\chapter{Canonical Form Reduction}\label{ch:canonical}

This is the first of three chapters that together prove the Fundamental Theorem
of matrix product states. Here an arbitrary translation-invariant MPS tensor is
reduced to a weighted family of primitive blocks; Chapter~\ref{ch:bnt} constructs
the basis of normal tensors carried by those blocks and compares their
coefficients, and Chapter~\ref{ch:ft_proof} proves the theorem itself. The
reduction follows~\cite{PerezGarcia2007Matrix}, \cite{Cirac2016MPDO_arXiv},
\cite{Cirac2021Matrix}, and~\cite[Chapter~6]{Wolf2012Quantum}. Its goal is to expose the primitive blocks
inside the tensor. The transfer operator of an arbitrary tensor can be reducible and periodic.
The reduction first splits invariant subspaces. For each irreducible
trace-preserving block, it then powers the transfer map, extracts the
adjoint-fixed cyclic projections, compresses to their corners, and proves that
each corner is primitive and tensor-irreducible. The blocks are returned in a
fixed gauge. The periodic decomposition theorem
of~\cite[Theorem~5]{PerezGarcia2007Matrix} gives $p$ translated
$p$-periodic state summands when $p\mid N$ and a zero state otherwise. We state
this fixed-length decomposition below before using its blocked cyclic-sector
counterpart in Theorem~\ref{thm:cyclic_sector_decomp_after_blocking}.

\begin{definition}[Periodic projector-component chain]
    \label{def:pgvwc07_projector_component_chain}
    \lean{MPSChainTensor.repeatPeriodPattern,
        MPSTensor.projectorComponentPattern,
        MPSTensor.projectorComponentChain}
    \leanok
    Let $A=(A^i)_{i=0}^{d-1}$ have bond dimension $D$, let
    $(Q_u)_{u\in\mathbb Z/p\mathbb Z}$ be cyclic projections, and fix a starting
    sector $u$. For each $0\le j<N$, the corresponding chain of length $N$
    has local matrices
    \begin{align}
        B_{u,j}^i
        &=Q_{u+j}A^iQ_{u+j+1},
        \label{eq:pgvwc07_projector_component_chain}
    \end{align}
    where sector indices are read modulo $p$. It has the same bond dimension
    $D$ as $A$ and repeats with period $p$.
\end{definition}

% Scope restriction (normalized orientation): this is the IsPeriodic surface,
% not yet a wrapper from PGVWC07's one-block unital hypotheses.
\begin{theorem}[Normalized periodic state-vector decomposition]
    \label{thm:pgvwc07_periodic_state_vector_decomposition}
    \lean{MPSTensor.pgvwc07_periodic_stateVector_decomposition_of_dvd,
        MPSTensor.pgvwc07_stateVector_eq_zero_of_not_dvd}
    \leanok
    \uses{def:pgvwc07_projector_component_chain, def:periodic_tensor}
    This is a normalized form of the state-vector conclusion in
    \cite[Theorem~5]{PerezGarcia2007Matrix}. Let $A$ be an irreducible
    left-canonical tensor whose peripheral transfer
    spectrum is the set of $p$-th roots of unity. For a positive ring length
    $N$, if $p\mid N$, then its state vector is the sum of the $p$ component
    chains in~(\ref{eq:pgvwc07_projector_component_chain}):
    \begin{align}
        V_N(A)
        &=\sum_{u\in\mathbb Z/p\mathbb Z}c_{B_u}.
        \label{eq:pgvwc07_periodic_state_vector_decomposition}
    \end{align}
    Each summand is represented by a $p$-periodic MPS of bond dimension $D$.
    If $p\nmid N$, then $V_N(A)=0$.
\end{theorem}

\begin{proof}
    \leanok
    \uses{def:pgvwc07_projector_component_chain}
    Choose orthogonal cyclic projections with
    $\sum_uQ_u=\Id$ and $Q_uA^i=A^iQ_{u+1}$. The component product telescopes to
    \begin{align}
        B_{u,0}^{i_0}\cdots B_{u,N-1}^{i_{N-1}}
        &=Q_uA^{i_0}\cdots A^{i_{N-1}}Q_{u+N}.
    \end{align}
    Summing its trace over $u$ gives the original coefficient. If $p\mid N$,
    then $Q_{u+N}=Q_u$ and every chain closes with period $p$. If $p\nmid N$,
    then $Q_{u+N}\ne Q_u$; cyclicity of the trace and orthogonality make every
    component coefficient zero.
\end{proof}

Two normalizations appear. The unital orientation
$\sum_i A_k^i(A_k^i)^\dagger=\Id$ is used in
\cite[Theorem~4]{PerezGarcia2007Matrix}; the trace-preserving
(left-canonical) orientation $\sum_i (A_k^i)^\dagger A_k^i=\Id$ is used in the
TP-gauge and after-blocking reductions below. They exchange under the
conjugate-transposed Kraus family $K_i:=(A^i)^\dagger$, whose transfer map is the
Frobenius adjoint of $\E_A$.

\section{Invariant-subspace decomposition and irreducible blocks}

The block decomposition rests on one principle: a positive semidefinite fixed
point of $\E_A$ that is not full rank has a proper invariant support, which
block-diagonalizes the Kraus operators without changing the MPV family. Iterating
the resulting splitting decomposes any tensor into irreducible blocks.

\begin{theorem}[Unitary conjugation preserves the MPV family]\label{thm:samempv_unitary}
    \lean{MPSTensor.sameMPV_conj_unitary}
    \leanok
    \uses{def:same_mpv}
    If $U$ is a unitary matrix, then $A$ and $U^\dagger A U$
    generate the same MPV family.
\end{theorem}

\begin{proof}\leanok
    By cyclicity of the trace,
    $\tr((U^\dagger A^{i_1} U) \cdots (U^\dagger A^{i_N} U))
        = \tr(A^{i_1} \cdots A^{i_N})$.
\end{proof}

\begin{theorem}[Invariant projection gives a two-block decomposition]\label{thm:lowerzero_two_block}
    \lean{MPSTensor.exists_twoBlock_decomp_of_lowerZero}
    \leanok
    \uses{def:same_mpv2}
    Suppose $P$ is an orthogonal projection satisfying, for every physical
    index~$i$,
    \[
        (\Id - P) A^i P = 0.
    \]
    Then there exist $n,m$ with $n + m = D$ and MPS tensors $A_1, A_2$ of bond
    dimensions $n$ and $m$ such that $A$ and the two-block tensor
    $A_1 \oplus A_2$ generate the same MPV family in the sense of
    Definition~\ref{def:same_mpv2}.
\end{theorem}

\begin{proof}
    \leanok
    \uses{thm:samempv_unitary}
    Diagonalize $P$ by a unitary $U$.
    Since $P^2 = P$, the diagonal form of $P$ has only $0$- and $1$-eigenvalues,
    so the basis splits as $n + m = D$.
    Conjugating $A$ by $U$ preserves the MPV family by
    Theorem~\ref{thm:samempv_unitary}.
    In this basis the relation $(\Id - P) A^i P = 0$ makes each letter upper
    block triangular.
    Write
    \[
        B^i =
        \begin{pmatrix}
            A_1^i & *     \\
            0     & A_2^i
        \end{pmatrix}
    \]
    for the conjugated letters in this split basis. Since
    $B^{i_1}\cdots B^{i_N}$ is again upper block triangular, its trace satisfies
    \[
        \tr(B^{i_1}\cdots B^{i_N})
        =
        \tr(A_1^{i_1}\cdots A_1^{i_N})
        +
        \tr(A_2^{i_1}\cdots A_2^{i_N}).
    \]
    Hence the block-diagonal tensor $A_1 \oplus A_2$ has the same MPV
    coefficients as the conjugated tensor.
\end{proof}

\begin{theorem}[Nontrivial invariant projection gives strict two-block reduction]
    \label{thm:lowerzero_two_block_strict}
    \lean{MPSTensor.exists_twoBlock_decomp_of_lowerZero_strict}
    \leanok
    \uses{def:same_mpv2}
    Under the hypotheses of Theorem~\ref{thm:lowerzero_two_block},
    assume in addition that $P \neq 0$ and $P \neq \Id$.
    Then there exist $n,m$ with $n + m = D$, $n < D$, and $m < D$, together
    with block tensors $A_1, A_2$, such that
    $A$ and $A_1 \oplus A_2$ generate the same MPV family in the sense of
    Definition~\ref{def:same_mpv2}.
\end{theorem}

\begin{proof}
    \leanok
    \uses{thm:lowerzero_two_block}
    Apply Theorem~\ref{thm:lowerzero_two_block}.
    In the diagonal form of $P$, the $1$-eigenspace is nonempty because
    $P \neq 0$, and the $0$-eigenspace is nonempty because $P \neq \Id$.
    Hence both block sizes are positive.
    Since $n + m = D$, this implies $n < D$ and $m < D$.
\end{proof}

\begin{definition}[Support projection]\label{def:support_proj}
    \lean{MPSTensor.supportProj}
    \leanok
    For a positive semidefinite matrix $\rho \ge 0$, the \emph{support
        projection} $P$ is the orthogonal projection onto the range of $\rho$. Via
    the spectral decomposition
    $\rho = U \diag(\lambda_1,\ldots,\lambda_D) U^\dagger$,
    \[
        P = U \diag(\mathbf{1}_{\lambda_1>0},\ldots,\mathbf{1}_{\lambda_D>0}) U^\dagger,
    \]
    where $\mathbf{1}_{\lambda_j>0}$ is $1$ if $\lambda_j > 0$ and $0$ otherwise.
\end{definition}

\begin{definition}[Invariant projection predicate]\label{def:has_inv_proj}
    \lean{MPSTensor.HasInvariantProj}
    \leanok
    \uses{def:mps_tensor}
    An MPS tensor $A$ \emph{has an invariant projection} if there exists an
    orthogonal projection $P$ with $P \neq 0$, $P \neq \Id$, and
    $(\Id - P) A^i P = 0$ for all~$i$. A tensor is \emph{irreducible} if it has
    no nontrivial invariant projection (Definition~\ref{def:irreducible_tensor}).
\end{definition}

\begin{theorem}[PSD fixed point gives invariant projection]\label{thm:fp_inv_proj}
    \lean{MPSTensor.lowerZero_of_posSemidef_fixedPoint}
    \leanok
    \uses{def:transfer_map}
    % Source: the support-projection step in the proof of Theorem 4
    % (PerezGarcia2007Matrix): the support of a singular positive fixed point is
    % invariant, A_i P_R = P_R A_i P_R, by positivity.
    Let $\rho \ge 0$ be a fixed point of $\E_A$, and let $P$ be its support
    projection. Then $(\Id - P) A^i P = 0$ for all~$i$: the range of $P$ is
    invariant under every $A^i$.
\end{theorem}

\begin{proof}
    \leanok
    From $\E_A(\rho) = \rho$ and positivity of each summand
    $A^i \rho (A^i)^\dagger$, one gets
    $(\Id - P) A^i \rho (A^i)^\dagger (\Id - P) = 0$ for each~$i$. Since $\rho$
    is supported on the range of $P$, this forces $(\Id - P) A^i P = 0$.

    This is the finite-dimensional case of~\cite[Proposition~6.10]{Wolf2012Quantum}:
    the support projection of any stationary state is sub-harmonic. The
    equivalence between irreducibility and absence of nontrivial invariant
    projections is~\cite[Theorem~6.2]{Wolf2012Quantum}; see
    also~\cite[Section~IV.A.1]{Cirac2021Matrix}.
\end{proof}

\begin{theorem}[Two-block decomposition]\label{thm:two_block_decomp}
    \lean{MPSTensor.exists_twoBlock_decomp_of_posSemidef_fixedPoint_strict}
    \leanok
    \uses{def:transfer_map, def:same_mpv2}
    % Source: the invariant-support projection and finite-ring trace split in
    % the proof of Theorem 4 (PerezGarcia2007Matrix).
    If $\E_A$ has a nonzero PSD fixed point $\rho$ that is not positive definite,
    then there exist MPS tensors $A_1, A_2$ of smaller bond dimensions with
    $\mc{V}(A) = \mc{V}(A_1 \oplus A_2)$.
\end{theorem}

\begin{proof}
    \leanok
    \uses{thm:fp_inv_proj, thm:samempv_unitary}
    The support projection $P$ of $\rho$ is neither $0$ nor $\Id$ (since
    $\rho \neq 0$ and $\rho$ is not positive definite).
    Theorem~\ref{thm:fp_inv_proj} gives $(\Id - P) A^i P = 0$, so a unitary
    change of basis block-diagonalizing $P$ puts $A^i$ in upper-triangular block
    form with diagonal blocks $B_{11}^i$, $B_{22}^i$.
    Theorem~\ref{thm:samempv_unitary} preserves the MPV under the conjugation,
    and the trace of an upper-triangular product equals the sum of the
    diagonal-block traces, so deleting the off-diagonal blocks preserves all MPV
    coefficients.
\end{proof}

\subsection{Iterated reduction to irreducible blocks}

Iterating the two-block splitting on bond dimension decomposes any tensor into
irreducible blocks; the remaining lemmas record the spectral facts about
irreducible unital blocks and the algebra-spanning input used later.

\begin{definition}[Irreducible tensor]\label{def:irreducible_tensor}
    \lean{MPSTensor.IsIrreducibleTensor}
    \leanok
    \uses{def:has_inv_proj}
    An MPS tensor~$A$ is \emph{irreducible} if it has no nontrivial invariant
    projection.
\end{definition}

\begin{theorem}[Irreducible block decomposition]\label{thm:irred_decomp}
    \lean{MPSTensor.exists_irreducible_blockDecomp}
    \leanok
    \uses{def:irreducible_tensor}
    Every MPS tensor $A$ admits a block-diagonal tensor $\bigoplus_{k=1}^r A_k$
    with each $A_k$ irreducible,
    $\mc{V}(A) = \mc{V}(\bigoplus_k A_k)$, and
    $\sum_{k=1}^rD_k=D$.
\end{theorem}

\begin{proof}
    \leanok
    \uses{thm:two_block_decomp}
    By strong induction on $D$. If $A$ is irreducible, take $r = 1$. Otherwise a
    nontrivial invariant projection $P$ gives the upper-triangular splitting of
    Theorem~\ref{thm:two_block_decomp}; deleting the off-diagonal part preserves
    the MPV family and yields two blocks of strictly smaller bond dimension.
    Apply the induction hypothesis to each block.

    Both~\cite[Theorem~4]{PerezGarcia2007Matrix}
    and~\cite[Section~IV.A.1]{Cirac2021Matrix} obtain the block decomposition by
    the same invariant-projection splitting.
\end{proof}

\section{Nonzero blocks and the trace-preserving gauge}%
\label{sec:nonzero_blocks_tp_gauge}

In an irreducible block decomposition some blocks may have all-zero Kraus
operators. Such a block contributes nothing to a matrix product vector at any
positive length and is therefore discarded. The structural direct-sum identity
implies that the retained block dimensions satisfy $\sum_k D_k\le D$ without
using an empty-word coefficient. Each retained block is then placed in the
trace-preserving gauge.

\begin{theorem}[Nonzero irreducible block decomposition]%
    \label{thm:nonzero_irreducible_block_decomposition}
    \lean{MPSTensor.exists_irreducible_blockDecomp_nonzeroBlocks}
    \leanok
    \uses{def:irreducible_tensor, def:same_mpv2_pos}
    Every MPS tensor $A$ admits a finite family of nonzero irreducible blocks
    $(B_k)_{k=1}^r$, each with at least one nonzero Kraus operator and positive
    bond dimension. For every $N \ge 1$ and every $\sigma$,
    \[
        V^{(N)}(A)_\sigma
        = V^{(N)}\!(\textstyle\bigoplus_{k=1}^r B_k)_\sigma.
    \]
    Moreover, $\sum_{k=1}^r D_k\le D$.
\end{theorem}

\begin{proof}\leanok
    \uses{thm:irred_decomp}
    Apply the irreducible block decomposition (Theorem~\ref{thm:irred_decomp})
    and retain the blocks having a nonzero Kraus operator. For every all-zero
    irreducible block $C_j$ and every $N \ge 1$, $V^{(N)}(C_j)_\sigma = 0$, so
    \[
        V^{(N)}(A)_\sigma
        = V^{(N)}\!(\textstyle\bigoplus_{k=1}^r B_k)_\sigma.
    \]
    The dimension inequality follows by restricting the structural identity
    $\sum_j D_j=D$ to the retained subset.
\end{proof}

\begin{theorem}[Blockwise TP-gauge normalization of irreducible blocks]%
    \label{thm:blockwise_tp_gauge}
    \lean{MPSTensor.exists_tp_gauge_blockwise}
    \leanok
    \uses{def:irreducible_tensor}
    Suppose $A$ generates the same MPV family as $\bigoplus_{k=1}^r B_k$, where
    each $B_k$ is irreducible with at least one nonzero Kraus operator. Then there
    are nonzero weights $(\mu_k)_{k=1}^r$ and blocks $(C_k)_{k=1}^r$ such that, for
    every $N$ and every $\sigma$,
    \[
        V^{(N)}(A)_\sigma
        =
        V^{(N)}\!(\textstyle\bigoplus_{k=1}^r \mu_k C_k)_\sigma
        =
        \sum_{k=1}^r \mu_k^{\,N} V^{(N)}(C_k)_\sigma,
    \]
    each $C_k$ irreducible and left-canonical,
    $\sum_i (C_k^i)^\dagger C_k^i = \Id$, with positive bond dimension. This is
    the blockwise form used in the arbitrary-tensor trace-preserving gauge
    below.
\end{theorem}

\begin{proof}\leanok
    The Perron--Frobenius TP-gauge construction is applied to each irreducible
    block. Gauge equivalence preserves irreducibility and the MPV family, and the
    positive scaling factors are absorbed into the weights $\mu_k$.
\end{proof}

\subsection{Trace-preserving gauge for arbitrary tensors}%
\label{sec:tp_gauge_arbitrary}

\begin{theorem}[TP-gauge reduction for arbitrary tensors]%
    \label{thm:tp_gauge_arbitrary}
    \lean{MPSTensor.exists_tp_gauge_from_arbitrary}
    \leanok
    \uses{def:irreducible_tensor, def:same_mpv2_pos}
    For any MPS tensor $A$, there exist nontrivial blocks $(B_k)_{k=1}^r$ with
    nonzero weights $(\mu_k)_{k=1}^r$,
    each block irreducible, left-canonical
    ($\sum_i (B_k^i)^\dagger B_k^i = \Id$), and of positive bond dimension, such
    that for every $N \ge 1$ and every $\sigma$,
    \[
        V^{(N)}(A)_\sigma
        = V^{(N)}\!(\textstyle\bigoplus_{k=1}^r \mu_k B_k)_\sigma,
    \]
    and $\sum_{k=1}^r D_k\le D$.
\end{theorem}

\begin{proof}\leanok
    \uses{thm:nonzero_irreducible_block_decomposition, thm:tp_data_irreducible}
    Theorem~\ref{thm:nonzero_irreducible_block_decomposition} supplies the nonzero
    irreducible blocks $(C_k)_{k=1}^r$. For every $N \ge 1$ and every $\sigma$,
    \[
        V^{(N)}(A)_\sigma
        = V^{(N)}\!(\textstyle\bigoplus_{k=1}^r C_k)_\sigma.
    \]
    The blockwise TP-gauge theorem then replaces each $C_k$ by a
    left-canonical block $B_k$ with weight $\mu_k$. Again for every $N \ge 1$
    and every $\sigma$,
    \[
        V^{(N)}\!(\textstyle\bigoplus_{k=1}^r C_k)_\sigma
        =
        V^{(N)}\!(\textstyle\bigoplus_{k=1}^r \mu_k B_k)_\sigma.
    \]
    The gauge preserves every block dimension, so the structural bound
    $\sum_kD_k\le D$ is unchanged.
\end{proof}

\section{Blocking, period removal, and primitive blocks}%
\label{sec:blocking_infrastructure}

Blocking by an integer~$p$ carries MPV equalities, weights, transfer maps, and
identifications of the blocked physical alphabet. The first consequence is the
existence of a common blocking period under which every transfer map in a finite
family becomes primitive; the period-removal and primitive-block reductions then
proceed from that common period.

\begin{theorem}[Common blocking period for a block family]%
    \label{thm:common_blocking_period}
    \lean{MPSTensor.exists_common_blocking_all_primitive}
    \leanok
    \uses{def:blocked_tensor}
    Let $(B_k)_{k=1}^r$ be a finite family of positive-bond-dimension MPS
    tensors, each admitting a blocking period $p_k$ making its transfer map
    primitive. Then $p = \lcm(p_1, \ldots, p_r)$ is a common
    period: $\E_{B_k^{[p]}}$ is primitive for every~$k$.
\end{theorem}

\begin{proof}\leanok
    Each $p_k$ divides $p$, so primitivity of $\E_{B_k^{[p_k]}}$ carries over to
    $\E_{B_k^{[p]}}$ via the divisibility identity for transfer-map powers.
\end{proof}

The following stability statements are used when a primitive irreducible block
is blocked again for a common-refinement or injectivity length. This later
blocking is distinct from the period-removal length that produced the sector
block.

\begin{theorem}[Blocking preserves irreducibility for primitive irreducible blocks]%
    \label{thm:blocked_irreducible_tp_primitive}
    \lean{MPSTensor.isIrreducibleTensor_blockTensor_of_tp_primitive_irr}
    \leanok
    \uses{def:blocked_tensor, def:irreducible_tensor}
    Let $A$ be left-canonical, irreducible, primitive, of positive bond
    dimension. For every $P>0$, the blocked tensor $A^{[P]}$ is irreducible:
    there is no orthogonal projection $Q$ with $Q \ne 0$, $Q \ne \Id$ such that,
    for every $P$-blocked index $\alpha$,
    \[
        (\Id-Q)(A^{[P]})^\alpha Q=0.
    \]
\end{theorem}

\begin{proof}\leanok
    The blocked transfer map is the corresponding positive power of the original
    one:
    \[
        \E_{A^{[P]}}(X)=\E_A^P(X).
    \]
    Hence a positive-definite fixed point of $\E_A$ remains fixed by
    $\E_{A^{[P]}}$. Primitivity gives uniqueness of positive semidefinite fixed
    points for the blocked map, hence its transfer map is irreducible, which is
    the irreducibility criterion for $A^{[P]}$.
\end{proof}

\begin{theorem}[Extra blocking of primitive irreducible sector blocks]%
    \label{thm:tp_primitive_irreducible_extra_blocking}
    \lean{MPSTensor.tp_primitive_irreducible_extra_blocking}
    \leanok
    \uses{def:blocked_tensor, thm:blocked_irreducible_tp_primitive}
    Let $A$ be trace-preserving, primitive, and tensor-irreducible. For every
    $k>0$,
    \[
        \sum_\alpha ((A^{[k]})^\alpha)^\dagger
        (A^{[k]})^\alpha = \Id,
        \qquad
        \sigma_\partial(\E_{A^{[k]}})=\{1\},
    \]
    and $A^{[k]}$ is tensor-irreducible. This $k$ is a later
    common-refinement or injectivity length, distinct from the period-removal
    length that produced the sector block.
\end{theorem}

\begin{proof}\leanok
    \uses{thm:blocked_irreducible_tp_primitive}
    Trace preservation follows from blocking a trace-preserving tensor,
    primitivity from a positive power of a primitive transfer map, and
    irreducibility from Theorem~\ref{thm:blocked_irreducible_tp_primitive}.
\end{proof}

\subsection{Period removal by cyclic sectors}%
\label{sec:cyclic_sectors}

The step needed below is the period-removing blocking described
in~\cite[Section~2.3]{Cirac2016MPDO_arXiv}: after blocking each irreducible TP
block by the least common multiple of its peripheral periods, the result has no
nontrivial $p$-periodic vectors. The cyclic-sector
decomposition comes from the peripheral spectrum of the adjoint transfer
map~\cite[Theorem~6.6]{Wolf2012Quantum}.

\begin{theorem}[Period removal by blocking]%
    \label{thm:cyclic_sector_decomp_after_blocking}
    \lean{MPSTensor.exists_cyclic_sector_decomp_after_blocking_of_TP_of_isIrreducibleTensor}
    \leanok
    \uses{def:blocked_tensor, thm:peripheral_cyclic_structure}
    Let $A$ be a left-canonical irreducible tensor. Then there is a
    period-removing length $m\ge 1$ such that $A^{[m]}$ admits:
    \begin{enumerate}
        \item left-canonical sector tensors $(C_k)_{k=1}^m$ of nonzero bond
              dimension,
        \item a unit-weight direct-sum MPV decomposition
              $A^{[m]} \sim \bigoplus_{k=1}^m C_k$,
        \item orthogonal projections $(P_k)_{k=1}^m$ summing to $\Id$ with
              $\E_{A^\dagger}(P_{k+1}) = P_k$,
        \item commutation relations
              $P_k A^{[m]}_i = A^{[m]}_i P_k$ for every blocked letter $i$,
        \item the sector trace formula
              $f_{C_k}(\sigma)=
                  \tr(P_k A^{[m]}_{\sigma_1}\cdots A^{[m]}_{\sigma_N})$,
        \item linear isomorphisms onto the compression corners
              $\varphi_k : \MN{\dim C_k} \xrightarrow{\sim} P_k \MN{D} P_k$,
        \item transfer-intertwining identities
              $\varphi_k(\E_{C_k^\dagger}(X)) =
                  \E_{(P_kA^{[m]})^\dagger}(\varphi_k(X))$,
        \item multiplicativity and adjoint compatibility
              $\varphi_k(XY)=\varphi_k(X)\varphi_k(Y)$ and
              $\varphi_k(X^\dagger)=\varphi_k(X)^\dagger$.
    \end{enumerate}
    This is the cyclic-sector part of the period-removal step
    of~\cite[Section~2.3]{Cirac2016MPDO_arXiv}; the projections $P_k$ come
    from the peripheral spectrum of the adjoint transfer map
    (\cite[Theorem~6.6]{Wolf2012Quantum}).
\end{theorem}

\begin{proof}\leanok
    The cyclic decomposition of the adjoint transfer map gives projections
    $(P_k)$ with $\E_{A^\dagger}(P_{k+1}) = P_k$. Iterating the cyclic relation
    $m$ times gives $(\E_{A^\dagger})^m(P_k) = P_k$, and the blocked-transfer
    identity $\E_{(A^{[m]})^\dagger} = (\E_{A^\dagger})^m$ identifies these with
    the adjoint transfer map of $A^{[m]}$. Apply the block decomposition from
    adjoint-fixed projections.
\end{proof}

\begin{theorem}[Unconditional primitive blocked sectors]%
    \label{thm:blocked_sector_unconditional}
    \lean{MPSTensor.primitive_and_irreducible_sectorBlocks_of_cyclic_decomp_after_blocking}
    \leanok
    \uses{thm:cyclic_sector_decomp_after_blocking,
        thm:sector_irred_unconditional,
        thm:primitive_adjoint_equiv}
    Let $A$ be a left-canonical irreducible tensor whose adjoint transfer map has
    peripheral spectrum $\{\gamma^j : 0 \le j < m\}$ for a primitive $m$th root of
    unity $\gamma$, with blocked cyclic-sector decomposition $(C_u)_{u=1}^m$,
    $(P_u)_{u=1}^m$, and
    linear isomorphisms
    $\varphi_u : \MN{\dim C_u} \xrightarrow{\sim} P_u \MN{D} P_u$. Assume
    explicitly that the $P_u$ are orthogonal projections summing to $\Id$ and
    satisfying $\E_{A^\dagger}(P_{u+1})=P_u$, that each sector has nonzero bond
    dimension, and that the compression maps satisfy
    \[
        \varphi_u(\E_{C_u^\dagger}(X))
        =
        \E_{(P_uA^{[m]})^\dagger}(\varphi_u(X)),
        \qquad
        \varphi_u(XY)=\varphi_u(X)\varphi_u(Y),
        \qquad
        \varphi_u(X^\dagger)=\varphi_u(X)^\dagger.
    \]
    Here $P_uA^{[m]}$ denotes the blocked tensor with letters
    $P_u(A^{[m]})^\alpha$.
    Then every
    sector tensor satisfies $\sigma_\partial(\E_{C_u})=\{1\}$ and is
    tensor-irreducible.
\end{theorem}

\begin{proof}\leanok
    \uses{thm:sector_irred_unconditional,
        thm:primitive_adjoint_equiv}
    Theorem~\ref{thm:sector_irred_unconditional} gives irreducibility of
    $(\E_A^\dagger)^m$ on each corner $P_u$. The compression equivalences
    $\varphi_u$ intertwine the adjoint transfer map of $C_u$ with the
    corresponding corner restriction of $(\E_A^\dagger)^m$. The peripheral
    spectrum of the blocked adjoint map is
    \[
        \sigma_\partial((\E_A^\dagger)^m)
        =
        \{(\gamma^j)^m : 0 \le j < m\}
        =
        \{1\},
    \]
    because $\gamma$ is a primitive $m$th root of unity. Hence each nonzero
    irreducible corner restriction is primitive. Primitivity and irreducibility
    are preserved across $\varphi_u$; then use
    Theorem~\ref{thm:primitive_adjoint_equiv} to pass from the adjoint blocked
    map to the primal transfer map of $C_u$.
\end{proof}

\begin{theorem}[Period removal with primitive irreducible sectors]%
    \label{thm:cyclic_sector_decomp_irr_tp_prim_irr}
    \lean{MPSTensor.exists_primitive_irreducible_cyclic_sector_decomp_of_TP_of_isIrreducibleTensor}
    \leanok
    \uses{thm:cyclic_sector_decomp_after_blocking, thm:blocked_sector_unconditional}
    Let $A$ be a left-canonical irreducible tensor. Then there is a period
    $m \ge 1$ and a unit-weight decomposition of $A^{[m]}$ into sector blocks
    $(C_k)_{k=1}^m$, each of which is left-canonical, has primitive transfer map,
    is tensor-irreducible, and has nonzero bond dimension. Equivalently, the
    cyclic-sector decomposition of $A^{[m]}$ already satisfies the
    trace-preservation, primitivity, and irreducibility hypotheses needed for the
    common-blocking step
    (Lemma~\ref{thm:exists_common_blocked_cyclic_sector_family}).
\end{theorem}

\begin{proof}\leanok
    \uses{thm:peripheral_cyclic_structure,
        thm:cyclic_sector_decomp_after_blocking, thm:blocked_sector_unconditional}
    Left-canonicality makes $K_i := (A^i)^\dagger$ unital, irreducibility passes
    to the adjoint map, and the quantum Perron--Frobenius theorem provides a
    positive-definite fixed point; the cyclic peripheral-spectrum theorem and
    Theorem~\ref{thm:cyclic_sector_decomp_after_blocking} then produce the
    cyclic-sector decomposition. Keep the projections and compression
    equivalences, and apply
    Theorem~\ref{thm:blocked_sector_unconditional} to the sector blocks.
\end{proof}

\begin{theorem}[Cyclic sectors after removing zero blocks]%
    \label{thm:ft_after_blocking_per_block_cyclic_nonzero_blocks}
    \lean{MPSTensor.afterBlocking_perBlockCyclicData_of_sameMPV₂Pos}
    \leanok
    \uses{thm:tp_gauge_arbitrary, def:same_mpv2_pos,
        def:primitive_irreducible_cyclic_sectors,
        thm:cyclic_sector_decomp_irr_tp_prim_irr}
    If $A$ and $B$ generate the same MPV family at every positive length, then after removing the
    all-zero blocks and applying the trace-preserving gauge there are weighted nonzero-block tensors
    $\bigoplus_j \mu^A_j A_j$ and $\bigoplus_k \mu^B_k B_k$ with, for every
    $N\ge 1$ and $\sigma$,
    \[
        V^{(N)}(A)_\sigma
        =
        V^{(N)}\!(\textstyle\bigoplus_j \mu^A_j A_j)_\sigma,
        \qquad
        V^{(N)}(B)_\sigma
        =
        V^{(N)}\!(\textstyle\bigoplus_k \mu^B_k B_k)_\sigma,
    \]
    and the two nonzero parts agree:
    \[
        V^{(N)}\!(\textstyle\bigoplus_j \mu^A_j A_j)_\sigma
        =
        V^{(N)}\!(\textstyle\bigoplus_k \mu^B_k B_k)_\sigma.
    \]
    Each nonzero block has primitive irreducible cyclic sectors: for each $j$,
    there is $m_j\ge 1$ with a unit-weight sector decomposition
    $A_j^{[m_j]} \sim \bigoplus_{u=1}^{m_j} C^A_{j,u}$, and similarly for $B_k$.
\end{theorem}

\begin{proof}\leanok
    \uses{thm:tp_gauge_arbitrary,
        thm:cyclic_sector_decomp_irr_tp_prim_irr}
    Apply Theorem~\ref{thm:tp_gauge_arbitrary} separately to $A$ and $B$. Its
    positive-length conclusion gives, for $N\ge 1$,
    \[
        V^{(N)}(A)_\sigma
        =
        V^{(N)}\!(\textstyle\bigoplus_j \mu^A_j A_j)_\sigma,
        \qquad
        V^{(N)}(B)_\sigma
        =
        V^{(N)}\!(\textstyle\bigoplus_k \mu^B_k B_k)_\sigma.
    \]
    Chaining these through $V^+(A)=V^+(B)$ identifies the two nonzero parts. Finally
    apply Theorem~\ref{thm:cyclic_sector_decomp_irr_tp_prim_irr} to
    every trace-preserving irreducible nonzero-weight block.
\end{proof}

\begin{theorem}[Common blocking length for cyclic sectors]%
    \label{thm:ft_after_blocking_common_length_common_sector_theorem}
    \lean{MPSTensor.afterBlocking_commonLengthCommonSectorData_of_sameMPV₂Pos}
    \leanok
    \uses{thm:ft_after_blocking_per_block_cyclic_nonzero_blocks, def:same_mpv2_pos,
        thm:exists_common_blocked_cyclic_sector_family_common_multiple,
        thm:same_mpv2_pos_block_tensor_to_tensor_from_blocks,
        thm:same_mpv2_pos_to_tensor_from_blocks_block_power,
        thm:common_blocked_cyclic_sector_derived_properties}
    If two tensors generate the same MPV family at every positive length, then the cyclic-sector
    decompositions on the two sides can be expressed at one common positive
    blocking length $p$. For every $N\ge 1$ and every blocked $\sigma$, each
    blocked tensor equals its powered nonzero part,
    \[
        V^{(N)}\!(A^{[p]})_\sigma
        =
        V^{(N)}\!(\textstyle\bigoplus_j (\mu^A_j)^p (A_j)^{[p]})_\sigma,
        \qquad
        V^{(N)}\!(B^{[p]})_\sigma
        =
        V^{(N)}\!(\textstyle\bigoplus_k (\mu^B_k)^p (B_k)^{[p]})_\sigma,
    \]
    and the two powered nonzero parts agree,
    \[
        V^{(N)}\!\bigl(\textstyle\bigoplus_j
        (\mu^A_j)^p (A_j)^{[p]}\bigr)_\sigma
        =
        V^{(N)}\!\bigl(\textstyle\bigoplus_k
        (\mu^B_k)^p (B_k)^{[p]}\bigr)_\sigma.
    \]
\end{theorem}

\begin{proof}\leanok
    \uses{thm:ft_after_blocking_per_block_cyclic_nonzero_blocks,
        thm:exists_common_blocked_cyclic_sector_family_common_multiple,
        thm:same_mpv2_pos_block_tensor_to_tensor_from_blocks,
        thm:same_mpv2_pos_to_tensor_from_blocks_block_power,
        thm:common_blocked_cyclic_sector_derived_properties}
    Start from
    Theorem~\ref{thm:ft_after_blocking_per_block_cyclic_nonzero_blocks}. With
    $p_A,p_B$ the least common multiples of the period-removal lengths on the two
    sides, use $p=\lcm(p_A,p_B)$. The prescribed-common-multiple
    lemma constructs both one-sided sector families at this length, and the
    positive-length blocking lemmas carry the one-sided tensor/nonzero-part
    equalities over to the common alphabet; for every $N\ge 1$,
    \[
        V^{(N)}\!(A^{[p]})_\sigma
        =
        V^{(N)}\!\bigl(\textstyle\bigoplus_j
        (\mu^A_j)^p (A_j)^{[p]}\bigr)_\sigma.
    \]
    The same applies on the $B$ side, and the two-sided blocking lemma gives
    equality of the two powered nonzero parts.
\end{proof}

\section{Blocking and weighted direct sums}%
\label{subsec:internal_lean_only_labels}

Blocking by a length~$p$ commutes with weighted direct sums up to the $p$th power
of each weight: for nonzero weights $(\mu_k)$ and blocks $(B_k)$,
\[
    \left(\textstyle\bigoplus_k \mu_k B_k\right)^{[p]}
    \sim \textstyle\bigoplus_k \mu_k^{\,p}\, B_k^{[p]},
\]
where $\sim$ is equality of MPV families. When two block families are reblocked to
a common length $p = m_k e_k$, the $p$-blocked alphabet is identified with the
$e_k$-fold tensor power of the $m_k$-blocked alphabet, so each sector reblocked
through its own period agrees with the directly $p$-blocked block.

\begin{lemma}[Adjoint primitivity equivalence]%
    \label{thm:primitive_adjoint_equiv}
    \lean{MPSTensor.IsPrimitive.adjoint_iff}\leanok
    \uses{def:primitive_channel}
    Let $\E$ be a finite-dimensional completely positive map. Then $\E$ is
    primitive if and only if its Frobenius adjoint $\E^\dagger$ is primitive;
    equivalently, $\sigma_\partial(\E) = \{1\}$ iff
    $\sigma_\partial(\E^\dagger) = \{1\}$, since the peripheral spectrum of
    $\E^\dagger$ is the complex conjugate of that of $\E$.
\end{lemma}

\begin{lemma}[Primitivity under blocking by a multiple]%
    \label{thm:primitive_blocking_divisible}
    \lean{MPSTensor.isPrimitive_transferMap_blockTensor_of_dvd}\leanok
    \uses{def:blocked_tensor, def:primitive_channel}
    Let $A$ have positive bond dimension and let $p_0,p\ge 1$ with $p_0 \mid p$.
    If $\E_{A^{[p_0]}}$ is primitive, then $\E_{A^{[p]}}$ is primitive.
\end{lemma}

\begin{lemma}[Blocking the weighted block sum]%
    \label{thm:block_tensor_to_tensor_from_blocks}
    \lean{MPSTensor.sameMPV₂_blockTensor_toTensorFromBlocks}\leanok
    \uses{def:blocked_tensor, def:block_diagonal_tensor, def:same_mpv2}
    For nonzero weights $(\mu_k)_{k=1}^r$, blocks $(B_k)_{k=1}^r$, and any
    $p\ge 1$,
    \[
        (\textstyle\bigoplus_k \mu_k B_k)^{[p]}
        \sim  \textstyle\bigoplus_k \mu_k^{\,p}\, B_k^{[p]},
    \]
    where $\sim$ denotes generation of the same MPV family at every length.
\end{lemma}

\begin{lemma}[Positive-length equality survives blocking]%
    \label{thm:same_mpv2_pos_block_tensor}
    \lean{MPSTensor.sameMPV₂Pos_blockTensor}\leanok
    \uses{def:blocked_tensor, def:same_mpv2_pos}
    If two tensors have the same MPV coefficients at every positive length, then
    so do their blocked tensors after any positive blocking length.
\end{lemma}

\begin{lemma}[Positive-length blocking of a weighted block sum]%
    \label{thm:same_mpv2_pos_block_tensor_to_tensor_from_blocks}
    \lean{MPSTensor.sameMPV₂Pos_blockTensor_toTensorFromBlocks}\leanok
    \uses{def:blocked_tensor, def:block_diagonal_tensor, def:same_mpv2_pos,
        thm:same_mpv2_pos_block_tensor, thm:block_tensor_to_tensor_from_blocks}
    If $A$ and $\bigoplus_k \mu_k B_k$ have the same MPV coefficients at every
    positive length, then after any positive blocking length $p$, $A^{[p]}$ and
    $\bigoplus_k \mu_k^{\,p} B_k^{[p]}$ have the same MPV coefficients at every
    positive length.
\end{lemma}

\begin{lemma}[Positive-length equality of powered block sums]%
    \label{thm:same_mpv2_pos_to_tensor_from_blocks_block_power}
    \lean{MPSTensor.sameMPV₂Pos_toTensorFromBlocks_blockPower}\leanok
    \uses{def:blocked_tensor, def:block_diagonal_tensor, def:same_mpv2_pos,
        thm:same_mpv2_pos_block_tensor, thm:block_tensor_to_tensor_from_blocks}
    If $\bigoplus_j \mu^A_j A_j$ and $\bigoplus_k \mu^B_k B_k$ have the same MPV
    coefficients at every positive length, then after any positive common
    blocking length $p$, $\bigoplus_j (\mu^A_j)^p A_j^{[p]}$ and
    $\bigoplus_k (\mu^B_k)^p B_k^{[p]}$ have the same MPV coefficients at every
    positive length.
\end{lemma}

\begin{lemma}[Nonzero weights survive blocking]%
    \label{lem:block_weights_ne_zero}
    \lean{MPSTensor.blockWeights_ne_zero}\leanok
    If $\mu \in \C$ is nonzero and $p \ge 1$, then $\mu^p \neq 0$; blocking
    preserves the nonzero-weight hypothesis of a weighted block decomposition.
\end{lemma}

\begin{lemma}[Common reblocking of cyclic-sector families]%
    \label{thm:exists_common_blocked_cyclic_sector_family}
    \lean{MPSTensor.exists_commonBlockedCyclicSectorFamily_of_hasPrimitiveIrreducibleCyclicSectors}\leanok
    \uses{def:primitive_irreducible_cyclic_sectors, def:blocked_tensor}
    Given a finite family $(B_k)_{k=1}^r$, each with primitive irreducible cyclic
    sectors, there is a common blocking length $p$ such that each $B_k^{[p]}$
    admits a unit-weight direct-sum decomposition over a single common blocked
    alphabet $d^{[p]}$ into left-canonical blocks with primitive transfer maps,
    tensor irreducibility, and positive bond dimensions.
\end{lemma}

\begin{lemma}[Sector irreducibility for the adjoint blocked map]%
    \label{thm:sector_irred_unconditional}
    \lean{MPSTensor.isIrreducibleOnCorner_of_cyclic_decomp_mps}\leanok
    \uses{def:irreducible_tensor}
    Let $A$ be a left-canonical irreducible MPS tensor with $D>0$, let $m$ be the
    period of its cyclic-sector decomposition, and let $(P_u)_{u=1}^m$ be the
    cyclic projections of the adjoint transfer map with
    $\E_{A^\dagger}(P_{u+1}) = P_u$ and $\sum_u P_u = \Id$. Then the restriction
    of $(\E_{A^\dagger})^m$ to each corner $P_u \MN{D} P_u$ is irreducible.
\end{lemma}

\begin{lemma}[Iterated blocking at the transfer-map level]%
    \label{thm:iterated_blocking_transfer}
    \lean{MPSTensor.transferMap_blockTensor_mul}\leanok
    \uses{def:blocked_tensor, def:transfer_map}
    For any MPS tensor $A$ and $m,n \ge 0$,
    \[
        \E_{(A^{[m]})^{[n]}} =  \E_{A^{[mn]}}.
    \]
\end{lemma}

\begin{definition}[Primitive irreducible cyclic-sector decomposition]%
    \label{def:primitive_irreducible_cyclic_sectors}
    \lean{MPSTensor.HasPrimitiveIrreducibleCyclicSectors}\leanok
    \uses{def:blocked_tensor, def:irreducible_tensor, def:primitive_channel,
        def:tp_kraus}
    $A$ \emph{has primitive irreducible cyclic sectors} of period $m \ge 1$ if
    there is a family $(C_k)_{k=1}^m$ with $A^{[m]} \sim \bigoplus_{k=1}^{m} C_k$
    (unit weights), each $C_k$ left-canonical, with primitive transfer map,
    tensor-irreducible, and of positive bond dimension.
\end{definition}

\begin{lemma}[Common reblocking at a prescribed multiple]%
    \label{thm:exists_common_blocked_cyclic_sector_family_common_multiple}
    \lean{MPSTensor.exists_commonBlockedCyclicSectorFamily_of_commonMultiple}\leanok
    \uses{thm:exists_common_blocked_cyclic_sector_family}
    With hypotheses as in
    Lemma~\ref{thm:exists_common_blocked_cyclic_sector_family}, and given any
    common multiple $p$ of the per-block period-removal lengths $(m_k)$, the
    common reblocking can be performed at the prescribed length $p$.
\end{lemma}

\begin{lemma}[Structural properties of the common reblocked sectors]%
    \label{thm:common_blocked_cyclic_sector_derived_properties}
    \lean{MPSTensor.CommonBlockedCyclicSectorFamily.derived_properties}\leanok
    \uses{thm:exists_common_blocked_cyclic_sector_family}
    Each sector tensor in the common reblocked cyclic-sector family of
    Lemma~\ref{thm:exists_common_blocked_cyclic_sector_family} is
    trace-preserving ($\sum_i (B^i)^\dagger B^i = \Id$), has primitive transfer
    map ($\sigma_\partial(\E_B) = \{1\}$), is tensor-irreducible, and has positive
    bond dimension.
\end{lemma}

\section{Translation-invariant canonical form}

\begin{theorem}[Translation-invariant canonical form]%
    \label{thm:pgvwc07_ti_canonical_form}
    \lean{MPSTensor.exists_pgvwc07_normalized_exact_form_after_rescaling_allow_empty}
    \leanok
    % Source: PerezGarcia2007Matrix, Theorem 4.
    % Block decomposition, three block conditions, and bond-dimension bound.
    % Proof order to preserve: spectral-radius normalization and full-rank
    % fixed-point gauge; invariant support projection from a singular positive
    % fixed point; finite-ring trace split into supported/orthogonal summands;
    % iteration via a non-scalar fixed point; dual fixed point and final unitary
    % diagonalization. Derive the invariant support before any abstract
    % block-splitting lemma.
    Let $A$ be a translation-invariant MPS representation of bond
    dimension~$D$. There is a positive scalar $s$ and a finite family of
    blocks $A_k$, possibly empty, such that the globally rescaled tensor
    $s^{-1}A$ has the same MPV coefficients on every positive-length ring as
    $\bigoplus_k \mu_k A_k^i$.
    The weights are positive and satisfy $1\ge |\mu_k|$; if the family is
    nonempty, then some weight has norm~$1$. For each block,
    \begin{align}
        \sum_i A_k^i(A_k^i)^\dagger
        &= \Id,
        \notag\\
        \sum_i (A_k^i)^\dagger \Lambda_k A_k^i
        &= \Lambda_k.
        \label{eq:canonical_pgvwc_block_conditions}
    \end{align}
    Here $\Lambda_k$ is diagonal, positive, and full-rank, and each block has
    positive bond dimension. Moreover, the only fixed points of
    $\E_{A_k}(X)=\sum_i A_k^iX(A_k^i)^\dagger$ are the scalar multiples of
    $\Id$: for every $X$, the identity $\E_{A_k}(X)=X$ implies that there
    exists $c\in\C$ such that $X=c\Id$.
    The total bond dimension of the decomposition is at most~$D$.

    This is the positive-length form of
    \cite[Theorem~4]{PerezGarcia2007Matrix}, with the global
    spectral-radius normalization step made explicit.
\end{theorem}

\begin{remark}[Source proof order]\label{rem:pgvwc07_ti_canonical_source_order}
    The proof follows
    \cite[Theorem~4]{PerezGarcia2007Matrix} in order: normalize the
    spectral radius and use a full-rank positive fixed point to obtain the
    unital gauge; derive the invariant support projection from a singular
    positive fixed point by the positivity contradiction; split the finite-ring
    trace into the supported and orthogonal summands; iterate, using a
    non-scalar fixed point to split further until each block has only scalar
    fixed points; and finally repeat the fixed-point argument for the dual map
    and diagonalize the resulting full-rank positive fixed point by a unitary
    gauge. The primitive-block results below are used only after their
    hypotheses have been obtained from this argument.
\end{remark}

\begin{remark}[CPSV blocked canonical form (source statement)]%
    \label{rem:cpgsv_canonical_form_source}
    % Maintainer source anchors:
    % Papers/1606.00608/MPDO-22-12-17-2.tex: blocking removes periodic peripheral
    % components; normal-tensor and canonical-form definition eq:II_CF1 with the
    % weight normalization; any tensor becomes canonical after blocking. The
    % phase-class quotient used by the BNT supplier is prop:char-BNT, restated in
    % Appendix A. The two-layer expansion is eq:II_ABasicTensors/decBSV.
    % Papers/2011.12127/TN-Review-main.tex: invariant-subspace reduction; period
    % removal; def:4:normal-tensor-mps and eq:II_CF1; BNT
    % definition/characterization; two-layer BNT/copy display.
    The blocked canonical-form reduction of
    \cite[Section~2.3]{Cirac2016MPDO_arXiv} and
    \cite[Section~IV.A]{Cirac2021Matrix} states that an arbitrary
    translation-invariant MPS tensor becomes, after blocking by some period
    $p\ge 1$, the direct sum of an all-zero summand and a weighted family of
    normal tensors $\bigoplus_k \mu_k B_k$ with $|\mu_k|\le 1$ and at least one
    $|\mu_k|=1$. The closest proved results here are
    Theorem~\ref{thm:tp_primitive_blockdecomp} and
    Theorem~\ref{thm:paperbnt_supplier_after_blocking}; the source upper
    normalization $|\mu_k|\le 1$ with a unit-modulus witness is the one
    remaining input, recorded in
    Remark~\ref{rem:arbitrary_input_reduction_gap}.
\end{remark}

\section{Normal tensors and canonical forms}\label{sec:cpsv_canonical_defs}

This section records the CPSV definitions of \emph{normal tensor} and
\emph{canonical form} from~\cite{Cirac2016MPDO_arXiv}. The basis of normal
tensors is stated in Chapter~\ref{ch:bnt}. Write
\begin{align}
    \sigma_\partial(\E)
    &=
    \{\lambda\in\sigma(\E):|\lambda|=r(\E)\}
    \label{eq:canonical_peripheral_spectrum}
\end{align}
for the peripheral spectrum of a transfer map $\E$, where $r(\E)$ is its
spectral radius. The stronger predicates used later
(Definition~\ref{def:is_normal_canonical_form} and the canonical-form predicate
of Chapter~\ref{ch:bnt}) add left-canonical normalization and spectral
primitivity. The weight moduli are taken non-increasing, not strictly
decreasing; equal-modulus and repeated-copy sectors are retained and compared by
the coefficient argument of Chapter~\ref{ch:bnt}.

\begin{definition}[CPSV normal tensor (NT)]\label{def:nt_cpsv}
    \lean{MPSTensor.IsNormalTensor}
    \leanok
    A tensor $A$ of physical dimension~$d$ and bond dimension~$D$ is a
    \emph{normal tensor} if, after the spectral-radius normalization
    of~\cite{Cirac2016MPDO_arXiv}, (i)~$A$ admits no nontrivial invariant
    orthogonal projection and (ii)~the associated CPM
    $\E_A(X)=\sum_i A^iX(A^i)^\dagger$ has spectral radius
    $r(\E_A)=1$ and peripheral spectrum $\sigma_\partial(\E_A)=\{1\}$.
\end{definition}

\begin{remark}
    Condition~(i) cannot be dropped. A reducible tensor may have peripheral
    spectrum $\{1\}$ while still admitting a nontrivial invariant projection,
    so condition~(ii) alone does not characterize normality.
\end{remark}

The CPSV \emph{canonical form} (CF) of a tensor $A$ is the normal-block
decomposition $A^i\cong\bigoplus_{k=1}^r\mu_kA_k^i$, with weights normalized by
$|\mu_k|\le 1$ and with at least one $|\mu_k|=1$; see
\cite[Section~2.3]{Cirac2016MPDO_arXiv}. In the statements below, the structural
part is supplied by the stronger normal canonical form of
Definition~\ref{def:is_normal_canonical_form}, and the global weight
normalization is invoked as a separate source hypothesis where needed.

\begin{remark}[Relation to the strong predicates]
    A tensor $A$ which is irreducible
    (Definition~\ref{def:irreducible_tensor}), has transfer map of spectral
    radius one, and has primitive transfer map is normal in the sense of
    Definition~\ref{def:nt_cpsv}. The reverse implications from the stronger
    predicates of Definition~\ref{def:is_normal_canonical_form} and of
    Chapter~\ref{ch:bnt} require the fundamental-theorem step or the
    algebraic-to-spectral normality equivalence, established later.
\end{remark}

\subsection{Normalization conventions}\label{sec:transfer_map_normalization}

\begin{remark}[Scalar normalization conventions]
    The reductions use four scalar facts: scaling preserves injectivity; the
    transfer map under a scalar~$\zeta$ rescales by $|\zeta|^2$; unit-modulus
    phase rescaling preserves $\sum_i(A^i)^\dagger A^i=\Id$; and the MPV of a
    weighted block-diagonal tensor factorizes as
    $\mu_k^NV^{(N)}(A_k)_\sigma
    =\eta_k^N|\mu_k|^NV^{(N)}(A_k)_\sigma$, with
    $\eta_k:=\mu_k/|\mu_k|$.

    Left-canonical normalization is required before the blocks can be compared.
    Without it, the implication
    $\mc{V}(\bigoplus_k\mu_kA_k)=\mc{V}(\bigoplus_k\mu_kB_k)
    \Rightarrow\forall k,\ \mc{V}(A_k)=\mc{V}(B_k)$
    already fails for two scalar blocks: for $\mu=(1,2)$, $A=(1,3/2)$, and
    $B=(3,1/2)$, the two direct sums generate identical MPV families but the
    first blocks differ at $N=1$. Left-canonical normalization removes this
    particular ambiguity; the sectors are then matched, up to a permutation,
    by the exact fixed-length linear independence of Chapter~\ref{ch:bnt},
    which requires no ordering of the weight moduli.
\end{remark}

\section{Projector criteria and decomposition auxiliaries}

\begin{definition}[Invariant-projector closure]
    \label{def:invariant_projector_closure}
    \lean{MPSTensor.HasInvariantProjectorClosure}
    \leanok
    \uses{def:mps_tensor}
    An MPS tensor $A$ has \emph{invariant-projector closure} if every
    orthogonal projection $P$ satisfying $PA^i=PA^iP$ for all~$i$ also
    satisfies $A^iP=PA^iP$ for all~$i$.
    This is the projector hypothesis in the sufficient condition for canonical
    form of~\cite[Section~2]{Cirac2016MPDO_arXiv}.
\end{definition}

\begin{theorem}[Invariant projections are reducing under projector closure]
    \label{thm:invariant_projector_closure_commutes}
    \lean{MPSTensor.commutes_of_hasInvariantProjectorClosure_of_lowerZero}
    \leanok
    \uses{def:invariant_projector_closure}
    Suppose that $A$ has invariant-projector closure and that $P$ is an
    orthogonal projection satisfying $(\Id-P)A^iP=0$ for all~$i$. Then $P$
    reduces every tensor matrix: $PA^i=A^iP$ for all~$i$.
\end{theorem}

\begin{proof}
    \leanok
    \uses{def:invariant_projector_closure}
    The lower-corner relation gives
    $(\Id-P)A^i=(\Id-P)A^i(\Id-P)$. Applying projector closure to $\Id-P$
    gives $A^i(\Id-P)=(\Id-P)A^i(\Id-P)$, hence the upper corner
    $PA^i(\Id-P)$ also vanishes. Therefore $PA^i=PA^iP=A^iP$.
\end{proof}

\begin{theorem}[Invariant-projector closure gives a reducing projection]
    \label{thm:exists_reducing_projection_of_invariant_projector}
    \lean{MPSTensor.exists_reducing_projection_of_hasInvariantProj}
    \leanok
    \uses{def:has_inv_proj, def:invariant_projector_closure}
    If $A$ has invariant-projector closure and admits a nontrivial invariant
    projection, then it admits a nontrivial orthogonal projection commuting
    with every $A^i$.
\end{theorem}

\begin{proof}
    \leanok
    \uses{thm:invariant_projector_closure_commutes}
    Apply Theorem~\ref{thm:invariant_projector_closure_commutes} to the
    projection supplied by Definition~\ref{def:has_inv_proj}.
\end{proof}

\begin{lemma}[Support isometry of an orthogonal projection]
    \label{lem:projection_support_isometry}
    \lean{IsOrthogonalProjection.exists_support_isometry}
    \leanok
    \uses{def:orthogonal_projection_channel}
    Every orthogonal projection $P$ on $\C^D$ factors through an isometry:
    there are $n\in\N$ and $V\in\C^{D\times n}$ such that
    \begin{align}
        V^\dagger V
        &= \Id,
        \notag\\
        VV^\dagger
        &= P,
        \notag\\
        n
        &= \tr(P).
        \label{eq:canonical_projection_isometry}
    \end{align}
\end{lemma}

\begin{proof}
    \leanok
    Since $P$ is Hermitian and idempotent, its eigenvalues are $0$ and $1$.
    Writing $n=\tr(P)$, unitary diagonalization gives
    \begin{align}
        P
        &=
        U\diag(\underbrace{1,\ldots,1}_{n},0,\ldots,0)U^\dagger.
        \notag
    \end{align}
    The matrix $V\in\C^{D\times n}$ formed by the first $n$ columns of $U$
    satisfies (\ref{eq:canonical_projection_isometry}).
\end{proof}

\begin{lemma}[Corner tensors inherit projector closure]
    \label{lem:projector_closure_corner}
    \lean{MPSTensor.hasInvariantProjectorClosure_compress_of_commutes}
    \leanok
    \uses{def:invariant_projector_closure, def:orthogonal_projection_channel}
    Let $A$ have invariant-projector closure, and let $V$ be an isometry,
    $V^\dagger V=\Id$, whose range projection $VV^\dagger$ commutes with every
    $A^i$. Then the corner tensor $B^i=V^\dagger A^iV$ again has
    invariant-projector closure.
\end{lemma}

\begin{proof}
    \leanok
    \uses{def:invariant_projector_closure}
    Let $Q$ be an orthogonal projection with $QB^i=QB^iQ$ for all~$i$.
    Then $VQV^\dagger$ is an orthogonal projection. Commutation of
    $VV^\dagger$ with each $A^i$ gives $V^\dagger A^i=B^iV^\dagger$, so
    left-invariance of $Q$ yields
    \begin{align}
        (VQV^\dagger)A^i
        &=VQB^iV^\dagger
        =VQB^iQV^\dagger
        =(VQV^\dagger)A^i(VQV^\dagger).
        \notag
    \end{align}
    Projector closure for $A$ yields
    $A^i(VQV^\dagger)=(VQV^\dagger)A^i(VQV^\dagger)$, and compressing both
    sides between $V^\dagger$ and $V$ returns $B^iQ=QB^iQ$.
\end{proof}

\begin{lemma}[Matrix product vectors split along an isometric partition]
    \label{lem:mpv_partition_isometries}
    \lean{MPSTensor.sameMPV₂_toTensorFromBlocks_of_isometry_family}
    \leanok
    \uses{def:same_mpv2, def:block_diagonal_tensor}
    Let $V_1,\ldots,V_r$ be isometries, $V_k^\dagger V_k=\Id$, with
    $\sum_kV_kV_k^\dagger=\Id$, and suppose $A^iV_k=V_kB_k^i$ for some
    tensors $B_1,\ldots,B_r$. Then
    $\mc{V}(A)=\mc{V}(\bigoplus_{k=1}^rB_k)$.
\end{lemma}

\begin{proof}
    \leanok
    The intertwining relation propagates from letters to words, so
    $A^wV_k=V_kB_k^w$ for every word~$w$. Inserting the partition of unity
    $\sum_kV_kV_k^\dagger=\Id$ into the trace and cycling gives
    \begin{align}
        \tr(A^w)
        &=\sum_{k=1}^r\tr\!\left(V_k^\dagger A^wV_k\right)
        =\sum_{k=1}^r\tr(B_k^w),
        \notag
    \end{align}
    which is the matrix product vector of the direct sum.
\end{proof}

\begin{theorem}[Decomposition into irreducible corners under projector closure]
    \label{thm:projector_closure_irreducible_corners}
    \lean{MPSTensor.exists_irreducible_blockDecomp_with_isometry_of_hasInvariantProjectorClosure}
    \leanok
    \uses{def:invariant_projector_closure, def:irreducible_tensor,
        def:same_mpv2, def:block_diagonal_tensor}
    Suppose $A$ has invariant-projector closure. Then there are isometries
    $V_1,\ldots,V_r$ with $V_k^\dagger V_k=\Id$, pairwise orthogonal ranges,
    and $\sum_kV_kV_k^\dagger=\Id$, such that $A^iV_k=V_kA_k^i$ for the
    corner tensors $A_k^i=V_k^\dagger A^iV_k$, every corner tensor $A_k$ is
    irreducible, and $\mc{V}(A)=\mc{V}(\bigoplus_{k=1}^rA_k)$.
\end{theorem}

\begin{proof}
    \leanok
    \uses{def:has_inv_proj, thm:invariant_projector_closure_commutes,
        lem:projection_support_isometry, lem:projector_closure_corner,
        lem:mpv_partition_isometries}
    Proceed by strong induction on the bond dimension. If $A$ is irreducible,
    the single corner $V=\Id$ suffices. Otherwise, $A$ has a nontrivial
    invariant projection $P$, which commutes with every $A^i$ by
    Theorem~\ref{thm:invariant_projector_closure_commutes}. Factoring
    $P=V_1V_1^\dagger$ and $\Id-P=V_2V_2^\dagger$ through support isometries
    by Lemma~\ref{lem:projection_support_isometry} splits $A$ into two corner
    tensors of strictly smaller bond dimension. For $j=1,2$, commutation gives
    $A^iV_j=V_j(V_j^\dagger A^iV_j)$, and each corner inherits
    invariant-projector closure by Lemma~\ref{lem:projector_closure_corner}.
    The two recursively obtained families are combined. Ranges from different
    corners are orthogonal because
    \begin{align}
        V_1^\dagger V_2
        &=V_1^\dagger P(\Id-P)V_2
        =0.
        \notag
    \end{align}
    Equality of matrix product vectors follows from
    Lemma~\ref{lem:mpv_partition_isometries}.
\end{proof}

\begin{definition}[No nontrivial periodic vectors]
    \label{def:no_periodic_vectors}
    \lean{MPSTensor.HasNoPeriodicVectors}
    \leanok
    \uses{def:mps_tensor, def:transfer_map, def:irreducible_tensor}
    An MPS tensor $A$ has \emph{no nontrivial $p$-periodic vectors} if, for
    every isometry $V$ with $V^\dagger V=\Id$ and $A^iV=VB^i$ for all~$i$,
    with $B$ irreducible (equivalently, for every restriction $B$ of $A$ to a
    minimal invariant subspace), and for every $\rho>0$ and $r>0$ with
    $\E_B(\rho)=r\rho$, every eigenvalue of $\E_B$ of modulus~$r$ equals~$r$.
    The eigenvalue $r$ is the spectral radius of $\E_B$
    (Theorem~\ref{thm:spectral_radius_pf_irred}), so the condition states
    that no block of the invariant-subspace decomposition
    of~\cite[Section~2]{Cirac2016MPDO_arXiv}, rescaled to spectral radius
    one, has a peripheral eigenvalue $e^{2\pi iq/p}\neq1$. Eigenvectors with
    such eigenvalues are the $p$-periodic vectors
    of~\cite[Section~2]{Cirac2016MPDO_arXiv}.
\end{definition}

\begin{lemma}[Transfer map of an intertwined corner]
    \label{lem:transfer_map_corner_intertwine}
    \lean{MPSTensor.transferMap_conj_of_intertwine}
    \lean{MPSTensor.hasEigenvalue_transferMap_of_intertwine}
    \leanok
    \uses{def:transfer_map}
    Let $V$ satisfy $A^iV=VB^i$ for all~$i$. Then, for every $X$,
    \begin{align}
        \E_A(VXV^\dagger)
        &=V\E_B(X)V^\dagger.
        \label{eq:canonical_corner_transfer}
    \end{align}
    If moreover $V^\dagger V=\Id$, then every eigenvalue of $\E_B$ is an
    eigenvalue of $\E_A$.
\end{lemma}

\begin{proof}
    \leanok
    Expanding the left-hand side of (\ref{eq:canonical_corner_transfer})
    termwise, the intertwining $A^iV=VB^i$ and its adjoint
    $V^\dagger(A^i)^\dagger=(B^i)^\dagger V^\dagger$ move both factors past
    $V$ and $V^\dagger$. For the eigenvalue statement, if
    $\E_B(X)=\mu X$ with $X\neq0$, then
    $\E_A(VXV^\dagger)=\mu VXV^\dagger$, and $VXV^\dagger\neq0$ because
    compression with the isometry gives $V^\dagger(VXV^\dagger)V=X$.
\end{proof}

\begin{lemma}[Perron eigenpair of a nonzero irreducible tensor]
    \label{lem:irreducible_transfer_perron_pair}
    \lean{MPSTensor.exists_posDef_transferMap_eigenvector_of_irreducible}
    \leanok
    \uses{def:irreducible_tensor, def:transfer_map}
    Let $B$ be an irreducible MPS tensor of positive bond dimension, not all
    of whose matrices vanish. Then there are $\rho>0$ and $r>0$ with
    $\E_B(\rho)=r\rho$.
\end{lemma}

\begin{proof}
    \leanok
    \uses{thm:pf_eigenvector_existence, def:irreducible_cp}
    Irreducibility of the tensor makes the transfer map an irreducible
    completely positive map, and a nonzero matrix makes it nonzero, so it
    does not annihilate any nonzero PSD matrix. The Perron--Frobenius theorem
    for positive maps (Theorem~\ref{thm:pf_eigenvector_existence}) gives a
    nonzero PSD eigenvector with eigenvalue $r>0$, and irreducibility upgrades
    it to a positive-definite one
    (\cite[Theorem~6.3(2)]{Wolf2012Quantum}).
\end{proof}

\begin{lemma}[Spectral normalization of an irreducible block]
    \label{lem:irreducible_block_spectral_normalization}
    \lean{MPSTensor.isNormalTensor_invSqrt_smul_of_unique_peripheral}
    \leanok
    \uses{def:irreducible_tensor, def:transfer_map, def:nt_cpsv}
    Let $B$ be an irreducible MPS tensor of positive bond dimension with
    $\E_B(\rho)=r\rho$ for some $\rho>0$ and $r>0$, and suppose every
    eigenvalue of $\E_B$ of modulus~$r$ equals~$r$. Then $r^{-1/2}B$ is a
    normal tensor.
\end{lemma}

\begin{proof}
    \leanok
    \uses{def:nt_cpsv, thm:spectral_radius_pf_irred}
    The rescaling identity $\E_{r^{-1/2}B}(X)=r^{-1}\E_B(X)$ gives
    $\E_{r^{-1/2}B}(\rho)=r^{-1}r\rho=\rho$, so $\rho$ is a
    positive-definite fixed point of $\E_{r^{-1/2}B}$. The rescaled map
    remains irreducible and completely positive, so
    Theorem~\ref{thm:spectral_radius_pf_irred} identifies its spectral radius
    with the eigenvalue one. The eigenvalues of modulus one of
    $\E_{r^{-1/2}B}$ are the eigenvalues of $\E_B$ of modulus~$r$, divided by
    $r$; by hypothesis the only such eigenvalue is one, so the peripheral
    eigenvalue set is $\{1\}$. Irreducibility is unchanged by a nonzero
    rescaling.
\end{proof}

\begin{lemma}[Canonical inclusion of a dependent direct summand]
    \label{lem:matrix_sigma_block_inclusion}
    \lean{Matrix.sigmaBlockInclusion}
    \lean{Matrix.sigmaBlockInclusion_isometry}
    \lean{Matrix.blockDiagonal'_mul_sigmaBlockInclusion}
    \lean{Matrix.sigmaBlockInclusion_conjTranspose_mul_blockDiagonal'}
    \lean{Matrix.sigmaBlockInclusion_compression}
    \lean{Matrix.blockDiagonal'_eq_sum_sigmaBlockInclusion}
    \lean{Matrix.blockDiagonal'_commute_sigmaBlockInclusion_rangeProjection}
    \leanok
    Let $E_k:H_k\to\bigoplus_lH_l$ be the canonical inclusion and let
    $B=\bigoplus_lB_l$. Then
    \begin{align}
        E_k^\dagger E_k
        &=\Id,
        \notag\\
        BE_k
        &=E_kB_k,
        \notag\\
        E_k^\dagger B
        &=B_kE_k^\dagger,
        \notag\\
        E_k^\dagger BE_k
        &=B_k,
        \notag\\
        B
        &=\sum_kE_kB_kE_k^\dagger.
        \label{eq:canonical_direct_summand}
    \end{align}
    In particular, $B$ commutes with the range projection
    $E_kE_k^\dagger$.
\end{lemma}

\begin{proof}\leanok
    The identities in (\ref{eq:canonical_direct_summand}) follow directly
    from the defining block-diagonal entries. The relations
    $BE_k=E_kB_k$ and $E_k^\dagger B=B_kE_k^\dagger$ imply that $B$
    commutes with $E_kE_k^\dagger$.
\end{proof}

\begin{lemma}[Canonical hypotheses descend to dependent direct summands]
    \label{lem:canonical_hypotheses_dependent_direct_summand}
    \lean{MPSTensor.projectorClosure_and_noPeriodicVectors_block_of_reindex_eq_blockDiagonal}
    \leanok
    \uses{lem:matrix_sigma_block_inclusion,
        def:invariant_projector_closure, def:no_periodic_vectors}
    Suppose that, after a change of coordinates, $A$ is the dependent direct
    sum of tensors $B_k$. If $A$ has invariant-projector closure and no
    nontrivial periodic vectors, then every $B_k$ has both properties. The
    assertion also applies to summands of dimension zero.
\end{lemma}

\begin{proof}\leanok
    Transport the canonical inclusion of the $k$-th summand back to the bond
    coordinates of $A$. This is an isometry which intertwines $B_k$ with
    $A$ on both sides, and its range projection commutes with every letter of
    $A$. Projector closure therefore descends by exact compression. A
    periodic vector of $B_k$ would embed through the same isometry as a
    periodic vector of $A$.
\end{proof}

\begin{lemma}[Closed-chain equality from an exact isometric decomposition]
    \label{lem:same_mpv_exact_isometric_decomposition}
    \lean{MPSTensor.sameMPV₂Pos_toTensorFromBlocks_of_reconstruction}
    \leanok
    \uses{def:same_mpv2_pos, def:block_diagonal_tensor}
    Suppose that isometries $V_k$ satisfy
    \begin{align}
        V_k^\dagger A^i
        &=(\mu_kB_k^i)V_k^\dagger,
        \notag\\
        A^i
        &=\sum_kV_k(\mu_kB_k^i)V_k^\dagger.
        \label{eq:canonical_exact_reconstruction}
    \end{align}
    Then $A$ and $\bigoplus_k\mu_kB_k$ have the same matrix product vectors
    at every positive length.
\end{lemma}

\begin{proof}\leanok
    Expand the first letter by the reconstruction formula in
    (\ref{eq:canonical_exact_reconstruction}). In each summand, move the
    remaining word through $V_k^\dagger$ using the intertwining relation.
    Cyclicity of the trace places $V_k^\dagger V_k=\Id$ next to the weighted
    block word, leaving precisely its closed-chain trace.
\end{proof}

\begin{theorem}[Sufficient condition for canonical form]
    \label{thm:projector_closure_canonical_form}
    \lean{MPSTensor.%
        exists_normalTensor_blockDecomp_with_isometry_of_hasInvariantProjectorClosure}
    \lean{MPSTensor.exists_normalTensor_blockDecomp_sameMPV₂Pos_of_hasInvariantProjectorClosure}
    \leanok
    \uses{def:invariant_projector_closure, def:no_periodic_vectors,
        def:nt_cpsv, def:same_mpv2_pos, def:block_diagonal_tensor}
    Suppose $A$ has invariant-projector closure and no nontrivial
    $p$-periodic vectors. Then there are normal tensors
    $A_1,\ldots,A_r$ of positive bond dimensions $D_1,\ldots,D_r$, with
    $\sum_kD_k\le D$, and nonzero weights $\mu_1,\ldots,\mu_r$ such that
    \begin{align}
        \mc{V}^+(A)
        &=
        \mc{V}^+\!\left(\bigoplus_{k=1}^r\mu_kA_k\right).
        \label{eq:canonical_projector_decomp}
    \end{align}
    Moreover, there are isometries $V_k:\C^{D_k}\to\C^D$ with pairwise
    orthogonal ranges such that
    \begin{align}
        A^iV_k
        &=V_k(\mu_kA_k^i),
        \notag\\
        A^i
        &=\sum_kV_k(\mu_kA_k^i)V_k^\dagger.
        \label{eq:canonical_projector_isometries}
    \end{align}
    This is the sufficient condition for canonical form
    of~\cite[Section~2]{Cirac2016MPDO_arXiv}, with that paper's convention
    $\sum_kD_k\le D$ allowing zero blocks: a corner of $A$ carrying the
    zero tensor cannot be rescaled to spectral radius one and is omitted
    because it vanishes at every positive length.
\end{theorem}

\begin{proof}
    \leanok
    \uses{thm:projector_closure_irreducible_corners,
        lem:irreducible_transfer_perron_pair,
        lem:irreducible_block_spectral_normalization}
    Decompose $A$ into irreducible corners $B_k$ along isometries $V_k$
    by Theorem~\ref{thm:projector_closure_irreducible_corners}; every matrix
    product vector coefficient of $A$ is the sum of those of the corners. A
    corner whose matrices all vanish contributes nothing at positive length
    and is dropped. Every remaining corner has a Perron eigenpair
    $\E_{B_k}(\rho_k)=r_k\rho_k$ by
    Lemma~\ref{lem:irreducible_transfer_perron_pair}. The absence of
    $p$-periodic vectors, applied to the corner $B_k$, gives that every
    eigenvalue of $\E_{B_k}$ of modulus~$r_k$ equals~$r_k$, so
    $A_k:=r_k^{-1/2}B_k$ is normal by
    Lemma~\ref{lem:irreducible_block_spectral_normalization}. Setting
    $\mu_k:=r_k^{1/2}$ gives $\mu_kA_k=B_k$ at the level of matrices, so for
    every word $w$ of positive length~$N$ the weight cancels:
    \begin{align}
        \mu_k^N\tr(A_k^w)
        &=
        (r_k^{1/2})^N(r_k^{-1/2})^N\tr(B_k^w)
        =
        \tr(B_k^w).
        \notag
    \end{align}
    Summing over $k$ gives (\ref{eq:canonical_projector_decomp}), while the
    corner reconstruction gives (\ref{eq:canonical_projector_isometries}).
    The retained dimensions satisfy
    $\sum_kD_k\le\sum_k\tr(V_kV_k^\dagger)=\tr(\Id)=D$.
\end{proof}

\begin{theorem}[Non-scalar fixed point gives a further split]%
    \label{thm:unital_nonscalar_fixed_point_split}
    \lean{MPSTensor.exists_singular_posSemidef_fixedPoint_of_unital_nonScalar_fixedPoint}
    \lean{MPSTensor.exists_twoBlock_decomp_of_unital_nonScalar_fixedPoint}
    \leanok
    \uses{def:transfer_map, def:same_mpv2, thm:two_block_decomp}
    Suppose $A$ is unital, $\E_A(\Id)=\Id$. If $X$ is a Hermitian fixed point
    of $\E_A$ which is not a scalar multiple of $\Id$, then there is a nonzero
    PSD fixed point $\rho$ which is not positive definite, and consequently
    $A$ admits a two-block MPV-equivalent decomposition with both block
    dimensions strictly smaller than~$D$. This is the fixed-point shift and
    further decomposition step in
    \cite[Theorem~4]{PerezGarcia2007Matrix}.
\end{theorem}

\begin{proof}\leanok
    Let $\lambda_{\max}$ be the largest eigenvalue of $X$. Then
    $\rho:=\lambda_{\max}\Id-X$ is positive semidefinite, singular, and
    nonzero because $X$ is not scalar. Unitality and the fixed-point equation
    for $X$ give $\E_A(\rho)=\rho$. Apply
    Theorem~\ref{thm:two_block_decomp}. The non-scalar hypothesis is needed
    because scalar multiples of $\Id$ are fixed by every unital map and yield
    no nontrivial support projection.
\end{proof}

\section{Irreducible-block auxiliaries}

\begin{theorem}[Scalar fixed points in an irreducible unital block]%
    \label{thm:irreducible_unital_scalar_fixed_point}
    \lean{MPSTensor.fixed_eq_scalar_of_isIrreducibleTensor_unital}
    \leanok
    \uses{def:transfer_map, def:irreducible_tensor}
    If $A$ is irreducible and unital, $\E_A(\Id)=\Id$, and $\E_A(X)=X$, then $X$
    is a scalar multiple of $\Id$. This is the final fixed-point assertion in the
    proof of \cite[Theorem~4]{PerezGarcia2007Matrix}.
\end{theorem}

\begin{proof}\leanok
    Tensor irreducibility gives irreducibility of the completely positive
    transfer map. By Wolf's Theorem~6.6 for irreducible unital Kraus maps, every
    fixed point is scalar.
\end{proof}

\begin{theorem}[Unital gauge from a full-rank Perron--Frobenius eigenvector]%
    \label{thm:unital_gauge_of_pd_fixed_point}
    \lean{MPSTensor.spectralUnitalGauge_isUnital_of_transferMap_eigenvector}
    \lean{MPSTensor.unitalGauge_isUnital_of_transferMap_fixedPoint}
    \lean{MPSTensor.sameMPV_unitalGauge}
    \leanok
    \uses{def:transfer_map}
    Let $\rho$ be positive definite with $\E_A(\rho)=r\rho$ for some $r>0$. For
    $B^i := r^{-1/2}\rho^{-1/2}A^i\rho^{1/2}$,
    \begin{equation}\label{eq:canon_unital_gauge_identity}
        \sum_i B^i(B^i)^\dagger = \Id.
    \end{equation}
    If $r=1$, then the unscaled gauge $B^i := \rho^{-1/2}A^i\rho^{1/2}$ is unital
    and generates the same MPV family as $A$. This is the spectral-radius
    normalization and full-rank fixed-point gauge of
    \cite[Theorem~4]{PerezGarcia2007Matrix}.
\end{theorem}

\begin{proof}\leanok
    The positive square root of $\rho$ is invertible. Substituting
    $B^i=r^{-1/2}\rho^{-1/2}A^i\rho^{1/2}$ into $\sum_i B^i(B^i)^\dagger$ and
    using $\E_A(\rho)=r\rho$ gives
    (\ref{eq:canon_unital_gauge_identity}). For $r=1$ the change from $A$ to $B$
    is a similarity, so cyclicity of the trace gives equality of all MPV
    coefficients.
\end{proof}

\begin{lemma}[Range factorization of the support projection]
    \label{lem:support_proj_range_factorization}
    \lean{MPSTensor.exists_supportProj_eq_mul}
    \lean{MPSTensor.supportProj_mul_left_eq_self}
    \lean{MPSTensor.one_sub_supportProj_mul_mul_supportProj_eq_zero}
    \leanok
    \uses{def:support_proj}
    The support projection $\pi$ of a positive semidefinite matrix $\rho$
    factors as $\pi=\rho W$ for a spectral pseudo-inverse $W$, witnessing
    that the range of $\pi$ is the range of $\rho$.  For a rectangular
    matrix $Y$, the support projection $\pi$ of $YY^\dagger$ fixes $Y$,
    $\pi Y=Y$; and whenever $AY=YG$ for some matrix $G$ --- that is, $A$
    maps the range of $Y$ into itself --- the one-sided invariance
    $(\Id-\pi)A\pi=0$ holds.
\end{lemma}

\begin{proof}
    \leanok
    \uses{def:support_proj}
    In the spectral decomposition
    $\rho=U\diag(\lambda)U^\dagger$ take
    $W=U\diag(\lambda^{+})U^\dagger$, where $\lambda^{+}$ inverts the
    positive eigenvalues and vanishes on the kernel; then
    $\rho W=U\diag(\mathbf 1_{\lambda>0})U^\dagger=\pi$.  From
    $\pi(YY^\dagger)=YY^\dagger$ the matrix
    $((\Id-\pi)Y)((\Id-\pi)Y)^\dagger$ vanishes, hence $\pi Y=Y$.
    Finally $(\Id-\pi)A\pi=(\Id-\pi)AY(Y^\dagger W)
    =(\Id-\pi)Y\,G(Y^\dagger W)=0$.
\end{proof}

\begin{theorem}[Irreducibility under rescaled conjugation]
    \label{thm:irreducible_tensor_smul_conj}
    \lean{MPSTensor.isIrreducibleTensor_smul_conj}
    \leanok
    \uses{def:irreducible_tensor, def:support_proj,
        lem:support_proj_range_factorization}
    Let $B$ be an irreducible tensor, let $X$ be invertible, and let
    $c\ne 0$.  Then the rescaled conjugated tensor with letters
    $c\,X^{-1}B^vX$ is irreducible.  In particular the rescaled unital
    gauge of an irreducible tensor is irreducible; this is the
    normalization of \cite[lines~224--225]{Cirac2016MPDO_arXiv} applied to
    an irreducible corner.
\end{theorem}

\begin{proof}
    \leanok
    \uses{lem:support_proj_range_factorization}
    An invariant orthogonal projection $P\notin\{0,\Id\}$ of the
    conjugated tensor makes the columns of $Y=XP$ span a subspace invariant
    under every letter of $B$: $B^vY=YG_v$ with
    $G_v=X^{-1}B^vXP$.  The support projection of $YY^\dagger$ is then a
    nontrivial invariant orthogonal projection of $B$ by
    Lemma~\ref{lem:support_proj_range_factorization}: it is nonzero
    because $Y\ne 0$, and it is not the identity because any nonzero
    vector of the kernel of $P$ transports through $(X^\dagger)^{-1}$ to
    the kernel of $YY^\dagger$.  This contradicts irreducibility of $B$.
\end{proof}

\begin{theorem}[Unital Schwarz data for an irreducible corner]
    \label{thm:spectral_unital_gauge_schwarz_setup}
    \lean{MPSTensor.spectralUnitalGauge_schwarz_setup}
    \leanok
    \uses{def:irreducible_tensor, def:transfer_map,
        thm:unital_gauge_of_pd_fixed_point, thm:irreducible_tensor_smul_conj}
    Let $B$ be an irreducible tensor with a positive definite transfer
    eigenvector, $\E_B(\rho)=r\rho$ with $\rho>0$ and $r>0$.  Then the
    rescaled unital gauge $K^v=r^{-1/2}\rho^{-1/2}B^v\rho^{1/2}$ is a
    unital Kraus family, is an irreducible tensor with an irreducible
    transfer map, and its adjoint map has a positive definite fixed point.
    These are the hypotheses of the cyclic decomposition of the peripheral
    spectrum, \cite[Theorem~6.6]{Wolf2012Quantum}.
\end{theorem}

\begin{proof}
    \leanok
    \uses{thm:unital_gauge_of_pd_fixed_point, thm:irreducible_tensor_smul_conj,
        thm:psd_fp_pd_irred}
    Unitality is Theorem~\ref{thm:unital_gauge_of_pd_fixed_point};
    irreducibility is Theorem~\ref{thm:irreducible_tensor_smul_conj}.  The
    adjoint of a unital family is trace preserving, so its transfer map is
    a channel and has a positive semidefinite fixed point by the Ces\`aro
    argument; irreducibility of the adjoint map upgrades the fixed point to
    a positive definite one.
\end{proof}

\begin{lemma}[Eigenvalue transport to the unital gauge]
    \label{lem:eigenvalue_transport_unital_gauge}
    \lean{MPSTensor.hasEigenvalue_transferMap_spectralUnitalGauge}
    \leanok
    \uses{def:transfer_map}
    In the setting of
    Theorem~\ref{thm:spectral_unital_gauge_schwarz_setup}, every transfer
    eigenvalue $\mu$ of $B$ becomes the transfer eigenvalue $\mu/r$ of the
    rescaled unital gauge: the gauge conjugates the transfer map by the
    congruence with $\rho^{1/2}$ and divides it by $r$.  In particular the
    peripheral eigenvalues $e^{2\pi iq/p}$ of
    \cite[lines~225--230]{Cirac2016MPDO_arXiv} arise from the transfer
    eigenvalues of modulus $r$.
\end{lemma}

\begin{proof}
    \leanok
    If $\E_B(Z)=\mu Z$ then
    $Z'=\rho^{-1/2}Z\rho^{-1/2}$ is nonzero and satisfies
    $\E_K(Z')=(\mu/r)Z'$ by direct computation.
\end{proof}

\begin{theorem}[Unitary diagonal positive fixed point in TP gauge]\label{thm:form_II}
    \lean{MPSTensor.exists_unitary_diag_posDef_fixedPoint_of_TP_of_isIrreducibleTensor}
    \leanok
    \uses{def:irreducible_tensor, def:tp_kraus}
    Let $A$ be an irreducible MPS tensor with $\sum_i (A^i)^\dagger A^i = \Id$ and
    $D > 0$. Then there exist a unitary $U$ and a diagonal positive definite
    $\Lambda$ such that, for $B^i := U^\dagger A^i U$,
    \[
        \sum_i (B^i)^\dagger B^i = \Id,
        \qquad
        \E_B(\Lambda) = \Lambda.
    \]
\end{theorem}

\begin{proof}
    \leanok
    \uses{thm:psd_fp_pd_irred}
    The TP hypothesis makes $\E_A$ a channel with a nonzero PSD fixed point.
    Irreducibility makes $\E_A$ an irreducible CP map, and
    Theorem~\ref{thm:psd_fp_pd_irred} upgrades the fixed point to positive
    definite. The spectral theorem diagonalizes it by a unitary conjugation,
    which preserves the TP normalization.
\end{proof}

\begin{remark}[Burnside input for cumulative spanning]\leanok
    \uses{thm:irreducible_tensor_to_action, thm:burnside_matrix,
        lem:algspan_cumulative_span}
    The passage from irreducible tensors to full algebra spanning first passes
    through the irreducible action on $\C^D$. Burnside's theorem then says that
    the resulting irreducible subalgebra of $\MN{D}$ is all of $\MN{D}$. Once
    the algebra span is full, finite-dimensionality gives a single cumulative
    span that already equals $\MN{D}$.
\end{remark}

\begin{lemma}[The algebra span reaches a finite cumulative span]\label{lem:algspan_cumulative_span}
    \lean{MPSTensor.exists_cumulativeSpan_eq_top_of_algSpan_eq_top}
    \leanok
    \uses{def:alg_span, def:cumulative_span}
    If the generated unital algebra satisfies $\algspn(A) = \MN{D}$,
    then there exists $N$ such that $T_N(A) = \MN{D}$.
\end{lemma}

\begin{proof}
    \leanok
    Every element of $\algspn(A)$ lies in some cumulative span $T_n(A)$:
    the generators lie in $T_1(A)$, scalar multiples of the identity lie in
    $T_0(A)$, and the family $\{T_n(A)\}_{n \ge 0}$ is closed under addition
    and under multiplication with lengths adding.
    Since $\MN{D}$ is finite-dimensional, the ascending chain
    $T_0(A) \subseteq T_1(A) \subseteq \cdots$ stabilizes, so one can choose a
    single $N$ with $T_N(A) = \MN{D}$.
\end{proof}

\begin{lemma}[Irreducible action gives an irreducible tensor]\label{lem:irreducible_action_to_tensor}
    \lean{MPSTensor.isIrreducibleTensor_of_isIrreducibleAction}
    \leanok
    \uses{def:irreducible_action, def:irreducible_tensor}
    If the letters of $A$ act irreducibly on $\C^D$, then $A$ has no nontrivial
    invariant orthogonal projection.
\end{lemma}

\begin{proof}
    \leanok
    If $P$ were a nontrivial invariant orthogonal projection, then its range
    would be invariant under every letter of $A$.
    Irreducibility of the action forces this range to be either $0$ or all of
    $\C^D$, so $P$ must be $0$ or $\Id$, a contradiction.
\end{proof}

\section{Normal canonical form}\label{sec:normal_canonical_form}

This section fixes the two canonical-form predicates used downstream: the
block-injective form of Definition~\ref{def:is_canonical_form} and the spectral
\emph{normal} canonical form of Definition~\ref{def:is_normal_canonical_form}. It
then records how the self-overlap clause and gauge-phase equivalence follow from
primitivity, and closes with the periodicity-removal step that produces primitive
transfer maps.

Recall from Definition~\ref{def:mpv_overlap}
and Lemma~\ref{thm:overlap_trace} that the bilinear overlap is
\[
    O_{AB}(N) = \sum_{\sigma \in \{0,\ldots,d{-}1\}^N}
    V^{(N)}(A)_\sigma \overline{V^{(N)}(B)_\sigma} = \Tr(F_{AB}^N).
\]
This notation enters the self-overlap clause of the canonical-form predicate
below.

\begin{definition}[Canonical form conditions]\label{def:is_canonical_form}
    \lean{MPSTensor.IsCanonicalForm}
    \leanok
    \uses{def:injective}
    A collection of scaling factors $(\mu_k)_{k=1}^r$ and
    block tensors $(A_k)_{k=1}^r$ satisfies the
    \emph{canonical form conditions} if:
    \begin{enumerate}
        \item each $A_k$ is injective,
        \item each $A_k$ satisfies the TP normalization
              $\sum_i (A_k^i)^\dagger A_k^i = \Id$,
        \item the moduli are non-increasing:
              $|\mu_1| \ge |\mu_2| \ge \cdots \ge |\mu_r|$,
        \item all $\mu_k \neq 0$,
        \item every block bond dimension is positive: $D_k \ge 1$ for every~$k$,
        \item the self-overlap converges: $O_{A_k A_k}(N) \to 1$
              for each~$k$.
    \end{enumerate}
    Only this one-sided normalization is assumed;
    no unital condition is assumed here.
    Condition~(6) is a primitivity hypothesis:
    for a primitive channel
    (\cite[Theorem~6.7(3)]{Wolf2012Quantum}),
    $T^n(\rho) \to \rho_\infty$ for every initial state,
    which forces the self-overlap to converge to~$1$.
    See also~\cite[Theorem~6.8]{Wolf2012Quantum}
    for the completely-positive characterisation
    and Chapter~\ref{ch:spectral} for the transfer-operator gap
    proof.
\end{definition}

% rem:pgvwc07_vs_local_cf (deleted): Definition~\ref{def:is_canonical_form} is a
% local/conditional predicate assuming injectivity and the self-overlap limit;
% it is not the source canonical-form existence or uniqueness theorem of
% \cite[Theorem~4]{PerezGarcia2007Matrix}.

\begin{remark}[Normalization convention]\label{rem:tp_vs_unital}
    \cite[Theorem~4]{PerezGarcia2007Matrix} uses the unital gauge
    $\sum_i A_k^i (A_k^i)^\dagger = \Id$; here the blocks are in the dual TP gauge
    $\sum_i (A_k^i)^\dagger A_k^i = \Id$. A positive-definite fixed point of the
    adjoint transfer map provides the change of gauge between the two
    (Theorem~\ref{thm:left_canonical_gauge}).
\end{remark}

\begin{definition}[Normal canonical form conditions]\label{def:is_normal_canonical_form}
    \lean{MPSTensor.IsNormalCanonicalForm}\leanok
    \uses{def:irreducible_tensor, def:primitive_channel, def:tp_kraus}
    Scaling factors $(\mu_k)_{k=1}^r$ and block tensors $(A_k)_{k=1}^r$ satisfy
    the \emph{normal canonical form conditions} if:
    \begin{enumerate}
        \item each $A_k$ is irreducible,
        \item each $A_k$ satisfies $\sum_i (A_k^i)^\dagger A_k^i = \Id$,
        \item each $\E_{A_k}$ is primitive
              (equivalently, $\sigma_\partial(\E_{A_k}) = \{1\}$),
        \item $|\mu_1| \ge |\mu_2| \ge \cdots \ge |\mu_r|$,
        \item all $\mu_k \neq 0$,
        \item all bond dimensions are positive.
    \end{enumerate}
    Here ``normal'' is the spectral sense of~\cite{Cirac2016MPDO_arXiv}, which
    by~\cite[Proposition~3]{Sanz2010Quantum} is equivalent to the eventual
    block-injectivity formulation of Definition~\ref{def:normal}.
\end{definition}

% The weight moduli are non-increasing, not strictly decreasing; equal-modulus
% and repeated-copy sectors are compared by the sector-coefficient construction
% of Chapter~\ref{ch:bnt} (Definition~\ref{def:mpv_phase_class_data}, supplier
% Theorem~\ref{thm:paperbnt_supplier_prepared}), where a repeated sector
% contributes the power sum $c_N^{(j)} = \sum_q \mu_{j,q}^N$.

\begin{theorem}[Self-overlap is derived in normal canonical form]%
    \label{thm:ncf_overlap_tendsto_one}
    \lean{MPSTensor.IsNormalCanonicalForm.overlap_tendsto_one}\leanok
    \uses{def:is_normal_canonical_form, def:mpv_overlap}
    If $(\mu_k, A_k)$ satisfies the normal canonical form conditions, then
    $O_{A_k A_k}(N) \to 1$ for every block~$k$.
\end{theorem}

\begin{proof}\leanok
    \uses{def:is_normal_canonical_form, thm:primitive_compl_lt_one, thm:primitive_overlap}
    Fix a block~$k$. Irreducibility and the TP normalization give a nonzero
    positive semidefinite fixed point $\rho$ of $\E_{A_k}$ with fixed-point
    projector $P$; primitivity gives the complementary gap
    $\rho_{\spec}(\E_{A_k}-P) < 1$ (Theorem~\ref{thm:primitive_compl_lt_one}).
    Theorem~\ref{thm:primitive_overlap} then yields $O_{A_k A_k}(N) \to 1$.
\end{proof}

\begin{remark}
    See Theorem~\ref{thm:modulus_one_gauge_irr_tp} for the modulus-one
    eigenvalue rigidity statement for irreducible trace-preserving blocks.
\end{remark}

\begin{lemma}[Mixed-transfer spectral radius from unit-modulus overlap]%
    \label{lem:overlap_norm_one_implies_spectral_radius_ge_one}
    \lean{MPSTensor.mixedTransferSpectralRadius_ge_one_of_mpvOverlap_norm_tendsto_one}
    \leanok
    \uses{def:mpv_overlap}
    Let $A$ and $B$ have the same bond dimension. If
    $\|\braket{V^{(N)}(A)}{V^{(N)}(B)}\| \to 1$, then the mixed transfer map has
    spectral radius at least $1$.
\end{lemma}
\begin{proof}\leanok
    If the mixed transfer spectral radius were $<1$, the iterated mixed transfer
    map would tend to zero in operator norm, hence so would its trace; via
    $\Tr(F_{AB}^N) = \braket{V^{(N)}(A)}{V^{(N)}(B)}$ the overlap would tend
    to~$0$, contradicting the hypothesis.
\end{proof}

\begin{theorem}[Gauge-phase equivalence from unit-modulus overlap]%
    \label{thm:nt_overlap_norm_one_gauge}
    \lean{MPSTensor.gaugePhaseEquiv_of_overlap_norm_tendsto_one_of_irreducible_TP}
    \leanok
    \uses{def:irreducible_tensor, def:tp_kraus, def:gauge_phase_equiv,
        def:mpv_overlap}
    Let $A$ and $B$ be irreducible MPS tensors of the same bond dimension with
    $\sum_i (A^i)^\dagger A^i = \Id$ and $\sum_i (B^i)^\dagger B^i = \Id$. If
    $\|\braket{V^{(N)}(A)}{V^{(N)}(B)}\| \to 1$, then $A$ and~$B$ are gauge-phase
    equivalent.
    % Same-bond-dimension case of \cite[Lemma~equalMPS]{Cirac2016MPDO_arXiv};
    % proportionality of MPVs is not assumed.
\end{theorem}

\begin{proof}\leanok
    \uses{lem:overlap_norm_one_implies_spectral_radius_ge_one,
        thm:modulus_one_gauge_irr_tp}
    By Lemma~\ref{lem:overlap_norm_one_implies_spectral_radius_ge_one} the mixed
    transfer spectral radius is at least~$1$, and
    Theorem~\ref{thm:modulus_one_gauge_irr_tp} then gives gauge-phase equivalence.
\end{proof}

\subsection{Canonical form from primitivity}\label{sec:canonical_from_primitivity}

A normal block needs a primitive transfer map; an irreducible TP block may still
have peripheral spectrum larger than $\{1\}$, namely the roots of unity coming
from its period. Blocking by the least common multiple of those periods raises the
transfer map to a power whose peripheral spectrum is $\{1\}$, removing the
periodicity. In~\cite{Cirac2021Matrix} this is the passage from irreducibility to
primitivity by blocking to period one.
% rem:primitivity_reduces_assumptions / "two ways" (deleted): once
% peripheral-spectrum primitivity is known, the overlap-convergence hypothesis is
% redundant, derived from the transfer-map gap of Chapter~\ref{ch:spectral}; these
% are statements about blocks already in normal form, not the arbitrary-input
% existence theorem of \cite[Theorem~4]{PerezGarcia2007Matrix}.

\begin{theorem}[Blocking yields a peripheral-spectrum primitive transfer map]\label{thm:blocking_primitivity}
    \lean{MPSTensor.exists_blockTensor_isPrimitive_of_TP_of_isIrreducibleTensor}
    \leanok
    \uses{def:blocked_tensor, def:transfer_map, def:primitive_channel,
        def:tp_kraus, def:irreducible_tensor}
    Let $A$ be an irreducible MPS tensor with $D > 0$ and
    $\sum_i (A^i)^\dagger A^i = \Id$. Then there exists a blocking length $p > 0$
    such that $\sigma_\partial(\E_{A^{[p]}}) = \{1\}$.
\end{theorem}

\begin{proof}\leanok
    \uses{thm:peripheral_roots_irred_fp, thm:peripheral_finite,
        lem:common_power_eq_one, lem:peripheral_pow_singleton}
    Pass to the conjugate-transposed Kraus family $K_i := (A^i)^\dagger$, which is
    unital with irreducible Kraus map. The TP fixed-point theorem for $\E_A$
    provides a positive definite fixed point for the adjoint of $K$, so
    Theorem~\ref{thm:peripheral_roots_irred_fp} shows every peripheral eigenvalue
    is a root of unity. Since $\sigma_\partial(\E_A)$ is finite
    (Theorem~\ref{thm:peripheral_finite}),
    Lemma~\ref{lem:common_power_eq_one} gives a common exponent $p$ with
    $\mu^p = 1$ for all $\mu \in \sigma_\partial(\E_A)$, and
    Lemma~\ref{lem:peripheral_pow_singleton} gives
    $\sigma_\partial(\E_A^p) = \{1\}$. Blocking by $p$ takes the $p$th power of
    the transfer map, so $\E_{A^{[p]}}$ is primitive.
\end{proof}

\begin{remark}\label{rem:blocking_primitivity_appendixA}
    Theorem~\ref{thm:blocking_primitivity} is the periodicity-removal-by-blocking
    step of~\cite[Appendix~A]{Cirac2016MPDO_arXiv} for an irreducible block
    already in TP gauge.
\end{remark}

\section{Left-canonical and dual diagonalization of irreducible blocks}%
\label{sec:canonical_construction_1606}

Each irreducible block in trace-preserving gauge admits the PGVWC07 unital
orientation together with a diagonal positive-definite dual fixed point
(Theorem~\ref{thm:CFII_data}, the Perron--Frobenius gauge of
Chapter~\ref{ch:qpf}, and
Theorem~\ref{thm:irreducible_unital_scalar_fixed_point}); the full PGVWC07
canonical form is Theorem~\ref{thm:pgvwc07_ti_canonical_form}.

\begin{theorem}[Left-canonical normalization for irreducible TP blocks]%
    \label{thm:CFII_data}
    \lean{MPSTensor.exists_CFII_data_of_TP_of_isIrreducibleTensor}
    \leanok
    \uses{def:irreducible_tensor}
    Let $A$ be an irreducible MPS tensor in TP gauge
    ($\sum_i (A^i)^\dagger A^i = \Id$) with $D > 0$. Then there exist a unitary
    $U$ and a diagonal positive definite $\Lambda$ such that:
    \begin{enumerate}
        \item $B^i := U^\dagger A^i U$ is TP,
        \item $\E_B(\Lambda) = \Lambda$,
        \item $\Lambda$ is diagonal and positive definite,
        \item $A$ and $B$ generate the same MPV family.
    \end{enumerate}
    This is the left-canonical normalization step (Canonical Form~II)
    of~\cite[Appendix~A]{Cirac2016MPDO_arXiv}. The unitary conjugation is a
    second step inside TP gauge, performed only after TP normalization.
\end{theorem}

\begin{proof}\leanok
    \uses{thm:form_II, thm:samempv_unitary}
    Theorem~\ref{thm:form_II} provides the unitary and the diagonal PD fixed
    point, and unitary conjugation preserves trace-preservation and the MPV
    family (Theorem~\ref{thm:samempv_unitary}).
\end{proof}

\begin{theorem}[Dual diagonalization for irreducible unital blocks]%
    \label{thm:pgvwc07_dual_diag_unital}
    \lean{MPSTensor.exists_unitary_diag_posDef_adjointFixedPoint_of_unital_of_isIrreducibleTensor}
    \leanok
    \uses{def:irreducible_tensor}
    Let $A$ be an irreducible MPS tensor in unital gauge
    $\sum_i A^i(A^i)^\dagger=\Id$ with $D>0$. Then there exist a unitary $U$ and a
    diagonal positive definite $\Lambda$ such that, for $B^i:=U^\dagger A^iU$,
    \begin{equation}\label{eq:canon_dual_diag_unital_pair}
        \sum_i B^i(B^i)^\dagger=\Id,
        \qquad
        \sum_i (B^i)^\dagger\Lambda B^i=\Lambda,
    \end{equation}
    and $A$, $B$ generate the same MPV family. This is the dual-map fixed-point
    and unitary-diagonalization step of
    \cite[Theorem~4]{PerezGarcia2007Matrix}.
\end{theorem}

\begin{proof}\leanok
    \uses{thm:CFII_data, thm:samempv_unitary}
    Apply Theorem~\ref{thm:CFII_data} to the conjugate-transposed family
    $i\mapsto (A^i)^\dagger$: unitality of $A$ is trace preservation for this
    family, whose transfer map is the dual of $\E_A$. Translating back gives the
    equations~(\ref{eq:canon_dual_diag_unital_pair}).
\end{proof}

\begin{lemma}[Common scalar rescaling of block weights]%
    \label{lem:mpv_to_tensor_from_blocks_weight_mul_left}
    \lean{MPSTensor.mpv_toTensorFromBlocks_weight_mul_left}
    \leanok
    Multiplying every block weight in a block-diagonal MPS tensor by a scalar $c$
    multiplies every length-$N$ MPV coefficient by $c^N$.
\end{lemma}

\begin{proof}\leanok
    \uses{thm:mpv_decomposition}
    Expand the block-diagonal MPV coefficient as a finite sum over blocks and
    factor $c^N$ out of the sum.
\end{proof}

\begin{theorem}[Normal canonical form from a primitive block decomposition]%
    \label{thm:exists_normal_canonical_form}
    \lean{MPSTensor.exists_normalCanonicalForm_of_primitive_blockDecomp}
    \leanok
    \uses{def:is_normal_canonical_form, def:blocked_tensor}
    Suppose $A$ generates the same MPV family as $\bigoplus_k \mu_k A_k$, and that
    each $A_k$ is irreducible, satisfies $\sum_i (A_k^i)^\dagger A_k^i = \Id$, and
    has primitive transfer map, all $\mu_k \neq 0$, and all bond dimensions are
    positive. Then there exist a blocking length $p > 0$, weights $(\nu_k)$, and
    blocked tensors $(B_k)$ such that $A^{[p]}$ generates the same MPV family as
    $\bigoplus_k \nu_k B_k$, and $(\nu_k, B_k)$ satisfies the normal canonical
    form conditions. The weight moduli need not be distinct; equal-modulus blocks
    are allowed.
\end{theorem}

\begin{proof}\leanok
    \uses{def:is_normal_canonical_form}
    The given blocks are already primitive, so take $p = 1$. Reorder the blocks by
    non-increasing weight modulus, with ties allowed, and set $\nu_k := \mu_k$ and
    $B_k := A_k$ in that order. Irreducibility, TP normalization, primitivity,
    nonzero weights, and positive bond dimensions are preserved, so the reordered
    family satisfies Definition~\ref{def:is_normal_canonical_form} and generates
    the same MPV family as~$A$.
\end{proof}

\begin{remark}[Scope of Theorem~\ref{thm:exists_normal_canonical_form}]
    Theorem~\ref{thm:exists_normal_canonical_form} assumes all six block
    hypotheses in advance and produces the normal canonical form by reordering.
    It is not the arbitrary-input existence theorem
    of~\cite[Section~2.3]{Cirac2016MPDO_arXiv}; that primitive-block reduction,
    including its structural bond-dimension bound, is
    Theorem~\ref{thm:tp_primitive_blockdecomp} in
    Section~\ref{sec:reduction_to_normal_form}.
\end{remark}


\section{Support compression and period-removal auxiliaries}

\begin{theorem}[Supported corner compression with a support isometry]%
    \label{thm:supported_corner_compression_isometry}
    \lean{MPSTensor.exists_compressedTensor_of_supported_projection_with_letter_and_isometry}
    \leanok
    \uses{def:support_proj}
    Let $A$ be a tensor normalized on the support of an orthogonal projection
    $P$, so that $P A^i P=A^i$ and
    $\sum_i (A^i)^\dagger A^i=P$. Then there are a corner tensor $C$, a
    linear isomorphism
    \[
        \varphi:\MN{\dim C}\xrightarrow{\sim} P\MN{D}P,
    \]
    and a rectangular isometry $V:\C^{\dim C}\to\C^D$ such that
    \[
        V^\dagger V=\Id,\qquad VV^\dagger=P,\qquad
        \varphi(X)=VXV^\dagger.
    \]
    Moreover $\varphi(C^i)=A^i$ for every physical letter, and for every word
    $\sigma=(\sigma_1,\ldots,\sigma_N)$,
    \[
        V^{(N)}(C)_\sigma
        = \tr\!\bigl(P A^{\sigma_1}\cdots A^{\sigma_N}\bigr).
    \]
    \begin{proof}\leanok
        Diagonalize the projection $P$ by a unitary change of basis.  The
        inclusion of the support coordinates gives a rectangular isometry $V$
        with $V^\dagger V=\Id$ and $VV^\dagger=P$.  Compress the letters of
        $A$ to this support and let $\varphi$ be the corresponding corner
        isomorphism.  The support relation $P A^iP=A^i$ gives the
        corner-letter identity. For each word
        $\sigma=(\sigma_1,\ldots,\sigma_N)$ the MPV component is obtained from
        the trace computation
        \[
            V^{(N)}(C)_\sigma
            =\tr(C^{\sigma_1}\cdots C^{\sigma_N})
            =\tr(V^\dagger A^{\sigma_1}\cdots A^{\sigma_N}V)
            =\tr(PA^{\sigma_1}\cdots A^{\sigma_N}).
        \]
        The final equality uses cyclicity of the trace and $VV^\dagger=P$.
    \end{proof}
\end{theorem}

\begin{theorem}[Commuting projections give sector isometries]%
    \label{thm:commuting_projection_block_decomp_isometry}
    \lean{MPSTensor.exists_blockDecomp_of_commuting_projections_with_letter_and_isometry}
    \leanok
    \uses{thm:supported_corner_compression_isometry}
    Let a left-canonical tensor $A$ commute with a family of orthogonal
    projections $(P_k)$ summing to $\Id$. Then $A$ decomposes into compressed
    sector tensors $(C_k)$, with a unit-weight direct-sum MPV decomposition,
    corner isomorphisms $\varphi_k:\MN{\dim C_k}\to P_k\MN{D}P_k$, and
    rectangular isometries $V_k$ satisfying
    \[
        V_k^\dagger V_k=\Id,\qquad V_kV_k^\dagger=P_k,\qquad
        \varphi_k(X)=V_kXV_k^\dagger .
    \]
    The sector letters obey
    $\varphi_k(C_k^i)=P_k A^i P_k$.
    \begin{proof}\leanok
        Apply Theorem~\ref{thm:supported_corner_compression_isometry} to the
        compression of $A$ on each sector $P_k$.  Orthogonality and
        $\sum_k P_k=\Id$ give the direct-sum decomposition. For every word
        $\sigma=(\sigma_1,\ldots,\sigma_N)$ the partition of unity gives
        \[
            V^{(N)}(A)_\sigma
            =\tr\!\Bigl((\sum_k P_k)A^{\sigma_1}\cdots A^{\sigma_N}\Bigr)
            =\sum_k\tr(P_kA^{\sigma_1}\cdots A^{\sigma_N})
            =\sum_k V^{(N)}(C_k)_\sigma,
        \]
        which is the unit-weight matrix-product-vector identity.
    \end{proof}
\end{theorem}

\begin{theorem}[Adjoint-fixed projections give sector isometries]%
    \label{thm:adjoint_fixed_projection_block_decomp_isometry}
    \lean{MPSTensor.exists_blockDecomp_of_adjoint_fixed_projections_with_letter_and_isometry}
    \leanok
    \uses{thm:commuting_projection_block_decomp_isometry}
    Let $A$ be left-canonical and let $(P_k)$ be orthogonal projections
    summing to $\Id$ and fixed by the adjoint transfer map. Then the same
    sector decomposition and support isometries of
    Theorem~\ref{thm:commuting_projection_block_decomp_isometry} exist. In
    particular, the adjoint fixed-point condition $\E_A^\dagger(P_k)=P_k$
    gives
    \[
        P_k A^i=A^i P_k,
    \]
    which gives the corner-letter identity
    \[
        \varphi_k(C_k^i)=P_k A^i P_k.
    \]
    \begin{proof}\leanok
        \uses{thm:commuting_projection_block_decomp_isometry}
        The fixed-point equation $\E_A^\dagger(P_k)=P_k$, together with
        left-canonicality $\sum_i (A^i)^\dagger A^i=\Id$, places $P_k$ in the
        multiplicative domain of $\E_A^\dagger$, forcing
        \[
            P_k A^i=A^i P_k
        \]
        for every letter $A^i$.  The result then follows from
        Theorem~\ref{thm:commuting_projection_block_decomp_isometry}.
    \end{proof}
\end{theorem}

\begin{theorem}[Spectral cyclic sectors with support isometries]%
    \label{thm:cyclic_sector_decomp_after_blocking_isometry}
    \lean{MPSTensor.exists_cyclic_sector_decomp_after_blocking_with_letter_and_isometry}
    \leanok
    \uses{thm:peripheral_cyclic_structure,
        thm:adjoint_fixed_projection_block_decomp_isometry}
    Let $A$ be a left-canonical irreducible tensor, and suppose the adjoint
    transfer map $\E_{A^\dagger}$ has a positive-definite fixed point, is
    irreducible as a positive map, and has peripheral spectrum
    \[
        \sigma_\partial(\E_{A^\dagger})
        =
        \{\gamma^j : j\in \Z/m\Z\}
    \]
    for a primitive $m$th root of unity $\gamma$, with $m\ge 1$. Then the
    spectral cyclic-sector decomposition of $A^{[m]}$ has compression maps
    which may be chosen with rectangular support isometries $V_k$ satisfying
    \[
        V_k^\dagger V_k=\Id,\qquad V_kV_k^\dagger=P_k,\qquad
        \varphi_k(X)=V_kXV_k^\dagger .
    \]
    These isometries also realize the blocked corner letters
    $\varphi_k(C_k^\alpha)=P_k(A^{[m]})^\alpha P_k$.
    \begin{proof}\leanok
        \uses{thm:peripheral_cyclic_structure,
            thm:adjoint_fixed_projection_block_decomp_isometry}
        The cyclic peripheral-spectrum decomposition applied to
        $\E_{A^\dagger}$ provides projections $P_k$ satisfying
        $\E_{A^\dagger}(P_{k+1})=P_k$.  Iterating $m$ times gives
        $(\E_{A^\dagger})^m(P_k)=P_k$, and the identity
        $(\E_{A^\dagger})^m=\E_{(A^{[m]})^\dagger}$ shows that each $P_k$ is
        fixed by the adjoint transfer map of $A^{[m]}$. Applying
        Theorem~\ref{thm:adjoint_fixed_projection_block_decomp_isometry} to
        that blocked tensor gives the compressed sector tensors and the
        support isometries.
    \end{proof}
\end{theorem}

\subsection{Reduction to primitive blocks}%
\label{sec:reduction_to_normal_form}

This subsection completes the reduction after period removal. A
TP-primitive-irreducible block is normal, and an arbitrary tensor becomes, after a
positive blocking length, a weighted family of left-canonical primitive blocks
whose total bond dimension is bounded by the original bond dimension. For two
tensors with the same MPV family, the cyclic sectors can be expressed over one
common blocked alphabet, which is the form used by the Fundamental Theorem.

\begin{theorem}[TP-primitive irreducible tensor is normal]%
    \label{thm:normal_tp_primitive_irred}
    \lean{MPSTensor.isNormal_of_tp_primitive_irreducible}
    \leanok
    \uses{def:primitive_mps, def:irreducible_tensor, def:normal}
    Let $A$ be a left-canonical MPS tensor with $D > 0$, irreducible transfer
    map, and peripheral spectrum $\{1\}$. Then $A$ is normal.
\end{theorem}

\begin{proof}\leanok
    \uses{thm:normal_from_primitive_posdef, thm:psd_fp_pd_irred}
    Irreducibility upgrades the PSD fixed point to positive definite
    (Theorem~\ref{thm:psd_fp_pd_irred}), and
    Theorem~\ref{thm:normal_from_primitive_posdef} then gives normality.
\end{proof}

An arbitrary translation-invariant tensor reduces, after blocking, to a weighted
direct sum of primitive irreducible blocks.

\begin{theorem}[Arbitrary tensor to primitive block decomposition]%
    \label{thm:tp_primitive_blockdecomp}
    \lean{MPSTensor.exists_tp_primitive_blockDecomp_after_blocking}
    \leanok
    \uses{def:blocked_tensor, def:same_mpv2_pos}
    For any MPS tensor $A$, there exist $p\ge 1$, nonzero weights
    $(\mu_k)_{k=1}^r$, and blocks $(B_k)_{k=1}^r$ such
    that, for every blocked $\sigma$ of positive length $N$,
    \[
        V^{(N)}\!(A^{[p]})_\sigma
        = V^{(N)}\!(\textstyle\bigoplus_{k=1}^{r} \mu_k B_k)_\sigma.
    \]
    The retained block dimensions satisfy $\sum_{k=1}^{r}D_k\le D$. Each block
    is left-canonical and primitive,
    \[
        \sum_i (B_k^i)^\dagger B_k^i = \Id,
        \qquad
        \sigma_\partial(\E_{B_k})=\{1\},
    \]
    and each bond dimension is positive.
\end{theorem}

\begin{proof}\leanok
    \uses{thm:tp_gauge_arbitrary, thm:common_blocking_period,
        lem:block_weights_ne_zero,
        thm:same_mpv2_pos_block_tensor_to_tensor_from_blocks}
    Theorem~\ref{thm:tp_gauge_arbitrary} supplies the left-canonical nontrivial
    blocks, the positive-length equality, and the bound $\sum_kD_k\le D$.
    Theorem~\ref{thm:common_blocking_period} provides a common blocking period
    $p$ making all transfer maps primitive,
    Lemma~\ref{lem:block_weights_ne_zero} keeps the blocked weights nonzero, and
    Lemma~\ref{thm:same_mpv2_pos_block_tensor_to_tensor_from_blocks} carries
    the positive-length equality through blocking.
\end{proof}

\begin{theorem}[Unconditional common primitive block decompositions]%
    \label{thm:unconditional_common_primitive_irreducible_blocks}
    \lean{MPSTensor.unconditional_commonPrimitiveIrreducibleBlocks}
    \leanok
    \uses{thm:ft_after_blocking_common_length_common_sector_theorem,
        thm:common_blocked_cyclic_sector_reindexed_nonzero_part,
        def:same_mpv2_pos}
    Let $A$ and $B$ have the same MPV family at every positive length. Then there is a positive blocking
    length at which both blocked tensors have the same positive-length MPV
    family as weighted direct sums of trace-preserving, primitive,
    tensor-irreducible blocks with positive bond dimensions and nonzero weights.
    The two weighted nonzero-sector families have the same positive-length MPV
    family.
\end{theorem}

\begin{proof}
    \leanok
    \uses{thm:ft_after_blocking_common_length_common_sector_theorem,
        thm:common_blocked_cyclic_sector_reindexed_nonzero_part}
    Apply
    Theorem~\ref{thm:ft_after_blocking_common_length_common_sector_theorem} to
    choose the common blocking length and the two common cyclic-sector families.
    For either one of these families, write $G_k$ for its cyclic-sector blocks,
    and write $(\nu_x,C_x)$ for the weights and blocks of the flattened
    common-sector family from
    Lemma~\ref{thm:common_blocked_cyclic_sector_reindexed_nonzero_part}. The
    common-alphabet nonzero-part identity gives, for every length $N$ and
    blocked word $\sigma$,
    \[
        V^{(N)}\!(\textstyle\bigoplus_k \mu_k^p G_k^{[p]})_\sigma
        =
        V^{(N)}\!(\textstyle\bigoplus_x \nu_x C_x)_\sigma ,
    \]
    Applying this on the two sides of the comparison yields a positive length
    $p$ and weighted primitive block families such that, for every $N\ge 1$ and
    blocked word $\sigma$,
    \[
        V^{(N)}\!(A^{[p]})_\sigma
        =
        V^{(N)}\!(\textstyle\bigoplus_x \mu^A_x A_x)_\sigma,
        \qquad
        V^{(N)}\!(B^{[p]})_\sigma
        =
        V^{(N)}\!(\textstyle\bigoplus_y \mu^B_y B_y)_\sigma,
    \]
    with the two weighted nonzero parts equal at every positive length. The
    structural properties and the nonvanishing of the weights come from the
    common-sector families.
\end{proof}

A weighted block sum with distinct sectors is the simplest sector decomposition:
one block per sector, each of multiplicity one. Equal-modulus and repeated-copy
sectors are not single scalar powers, already for $A = C \oplus (-C)$ with
$V^{(N)}(A)_\sigma = (1 + (-1)^N)\,V^{(N)}(C)_\sigma$; they are compared by the
basis-of-normal-tensors construction of Chapter~\ref{ch:bnt}.

\begin{definition}[Single-copy sector decomposition]
    \label{def:trivial_sector_decomp}
    \lean{MPSTensor.trivialSectorDecomp}
    \leanok
    \uses{def:sector_weight_data}
    Given nonzero weights $(\mu_k)_{k=1}^r$ and blocks $(B_k)_{k=1}^r$, the
    \emph{single-copy sector decomposition} has $r$ sectors, one basis block
    $B_k$ per sector, multiplicity~$1$, and sector weight $\mu_k$. As a finite
    sum,
    \[
        \sum_{k=1}^{r} \mu_k^{\,N} V^{(N)}(B_k)_\sigma
        =  V^{(N)}\!(\textstyle\bigoplus_k \mu_k B_k)_\sigma.
    \]
\end{definition}

\begin{lemma}[Single-copy sector decomposition preserves the represented family]
    \label{lem:trivial_sector_decomp_same_mpv}
    \lean{MPSTensor.sameMPV₂_trivialSectorDecomp}
    \leanok
    \uses{def:trivial_sector_decomp}
    The single-copy sector decomposition of $(\mu_k,B_k)_k$ represents the same
    MPS family as $\bigoplus_k \mu_k B_k$.
\end{lemma}

\begin{proof}\leanok
    \uses{def:trivial_sector_decomp}
    Expanding the sector decomposition gives $\sum_k \mu_k^N V^{(N)}(B_k)_\sigma$,
    the finite block-sum expression for $\bigoplus_k \mu_k B_k$.
\end{proof}


\section{Further weighted-block identities}

\begin{lemma}[Canonical form from peripheral primitivity]%
    \label{thm:canonical_from_primitive}
    \lean{MPSTensor.IsCanonicalForm.of_peripheral_primitive}
    \leanok
    \uses{def:is_canonical_form, def:injective, def:tp_kraus,
        def:primitive_channel}
    Let $(\mu_k,A_k)_{k=1}^r$ be nonzero weights and injective left-canonical
    blocks with positive bond dimensions, non-increasing moduli
    $|\mu_1|\ge\cdots\ge|\mu_r|>0$, and $\sigma_\partial(\E_{A_k}) = \{1\}$. Then
    $(\mu_k,A_k)_{k=1}^r$ satisfies the canonical-form conditions of
    Definition~\ref{def:is_canonical_form}, and $O_{A_k A_k}(N) \to 1$ for every
    $k$.
\end{lemma}

\begin{proof}
    \leanok
    \uses{def:is_canonical_form, thm:primitive_compl_lt_one, thm:primitive_overlap}
    The injectivity, left-canonical equation, weight ordering, and nonzero-weight
    clauses are exactly the hypotheses.  For each block, peripheral primitivity
    gives the complementary spectral gap
    $\rho_{\spec}(\E_{A_k}-P) < 1$ by Theorem~\ref{thm:primitive_compl_lt_one}.
    Theorem~\ref{thm:primitive_overlap} then yields $O_{A_k A_k}(N)\to 1$.
\end{proof}

\begin{lemma}[Common-alphabet flattening]%
    \label{thm:common_blocked_cyclic_sector_reindexed_samempv}
    \lean{MPSTensor.CommonBlockedCyclicSectorFamily.reindexed_sameMPV₂}\leanok
    \uses{thm:exists_common_blocked_cyclic_sector_family}
    With $p = m_k e_k$ for each $k$, the iterated blocked tensor
    $(B_k^{[m_k]})^{[e_k]}$ agrees with $B_k^{[p]}$ in MPV family, after the
    canonical identification of the $p$-blocked alphabet $d^{[p]}$ with the
    $e_k$-fold tensor power of the $m_k$-blocked alphabet $d^{[m_k]}$.
\end{lemma}

\begin{lemma}[Reindexed nonzero-part identity]%
    \label{thm:common_blocked_cyclic_sector_reindexed_nonzero_part}
    \lean{MPSTensor.CommonBlockedCyclicSectorFamily.reindexed_nonzero_part}\leanok
    \uses{thm:common_blocked_cyclic_sector_reindexed_samempv,
        thm:common_blocked_cyclic_sector_block_tensor_common_reindexed}
    The powered nonzero-part decomposition
    $\bigoplus_k \mu_k^p B_k^{[p]}$ carries over under the common-alphabet
    identification of
    Lemma~\ref{thm:common_blocked_cyclic_sector_reindexed_samempv}: the
    common-alphabet sectors have the same nonzero-part decomposition as the direct
    $p$-blocked tensors.
\end{lemma}

\begin{lemma}[Direct blocking in the common alphabet]%
    \label{thm:common_blocked_cyclic_sector_block_tensor_common_reindexed}
    \lean{MPSTensor.CommonBlockedCyclicSectorFamily.blockTensor_sameMPV₂_commonReindexedBlock}\leanok
    For every original block $B_k$, direct blocking at the common length
    $p=m_k e_k$ agrees, as an MPV family, with the same block transported to the
    common blocked alphabet. Write $\widetilde B_k$ for this common-alphabet
    image. Then
    \[
        V^{(N)}\!(B_k^{[p]})_\sigma
        =
        V^{(N)}\!(\widetilde B_k)_\sigma.
    \]
\end{lemma}

